\documentclass[10pt,a4paper]{article}

% ----------------- Paquetes -------------------%
\usepackage[T1]{fontenc}
\usepackage[utf8]{inputenc}
\usepackage[a4paper,margin=2.5cm]{geometry}
\usepackage{mathptmx}             
\usepackage{changepage} 
\usepackage{setspace}
\usepackage[spanish]{babel}
\usepackage[labelfont=bf,labelsep=space]{caption}
\usepackage{amssymb}
\usepackage{amsmath}
\usepackage{ebgaramond}
\usepackage{placeins}
\usepackage{etoolbox}
\usepackage{setspace}
\usepackage{parskip} 
\usepackage{tcolorbox}
\usepackage{needspace}
\usepackage{xparse}
\usepackage{ocgx2}
\usepackage{hyperref}
\usepackage{hypcap}
\usepackage{soul}\sethlcolor{yellow}% resaltado amarillo: \hl{texto} (Panel TCEE, ⌃H)



%%%%%%%%%%%%%%%%%%%%%%%%%%%%%%%%%%%%%%%%%%%%%%%%%%%%%%%%%%%%%%%%
% --- Fija primero tus valores globales "buenos" ---
\setstretch{1.2}                 % Interlineado global
\setlength{\parskip}{0.7\baselineskip}
\setlength{\parindent}{0pt}
\newlength{\GlobalParskip}
\setlength{\GlobalParskip}{\parskip}

\newcounter{eqnum}[subsection]
\renewcommand{\theeqnum}{\arabic{eqnum}}
\newcounter{alphaitem}
\newcounter{numitem}

%% ------------------- Recuadros----------------------%%
\newcommand{\puntosclave}[1]{%
  \begin{tcolorbox}[
    colframe=black,
    title={Puntos clave},
    before upper={\setlength{\parindent}{0.5cm}\setlength{\parskip}{0.5\baselineskip}}
  ]
  #1
  \end{tcolorbox}
}

\newcommand{\modificaciones}[1]{
  \begin{tcolorbox}[
    colback=white,
    colframe=red,
    title={Modificaciones a realizar},
    before upper={\setlength{\parindent}{0.5cm}\setlength{\parskip}{0.5\baselineskip}}
  ]
  #1
  \end{tcolorbox}
}

%% ------------------- Preguntas TEST----------------------%%
% Opciones: radio (single-choice)
\NewDocumentCommand{\OptRadio}{m m m}{%
  \noindent\CheckBox[radio,name=#1,value=#2,width=1.2ex,height=1.2ex]{}%
  \;\textbf{#2.}~#3\par
}

% Opciones: checkbox (multi-choice)
\NewDocumentCommand{\OptCheck}{m m}{%
  \noindent\CheckBox[name=#1,width=1.2ex,height=1.2ex,bordercolor={0 0 0}]{}%
  \;#2\par
}

% Pregunta single-choice (radio)
% Args: 1=id, 2=enunciado, 3=opciones, 4=solución+justificación, 5=imagen (o {}), 6=origen
\NewDocumentCommand{\PreguntaTestSingle}{m m m m m m}{%
  \par\medskip
  \noindent\textbf{#1.}~#2\par\smallskip

    % Imagen (si no es {})
    \IfBlankTF{#5}{}{%
      \begin{center}
        \includegraphics[width=0.75\linewidth]{#5}
      \end{center}
    }

    \begin{Form}
      #3
    \end{Form}

    \par\smallskip
    \switchocg{sol-#1}{\fbox{Mostrar / ocultar solución}}%
    \hfill{\footnotesize #6}\par\smallskip

    \begin{ocg}{Solución}{sol-#1}{0}
      \noindent #4
    \end{ocg}
  \par\medskip
}

% Pregunta multi-choice (checkboxes)
\NewDocumentCommand{\PreguntaTestMulti}{m m m m m m}{%
  \par\medskip
  \noindent\textbf{#1.}~#2\par\smallskip

    % Imagen (si no es {})
    \IfBlankTF{#5}{}{%
      \begin{center}
        \includegraphics[width=0.75\linewidth]{#5}
      \end{center}
    }
    \begin{Form}
      #3
    \end{Form}

    \par\smallskip
    \switchocg{sol-#1}{\fbox{Mostrar / ocultar solución}}%
    \hfill{\footnotesize #6}\par\smallskip

    \begin{ocg}{Solución}{sol-#1}{0}
      \noindent #4
    \end{ocg}
  \par\medskip
}

%% ------------------- CITA ----------------------%%
% \begin{cita}[Autor][Año][Obra] texto \end{cita}: cita sangrada y, debajo a la derecha, «AUTOR (Año), Obra». Los tres datos son opcionales.
% En el Panel TCEE: escribir \cita e Intro (el cursor va al texto; Tab pasa a Autor, Año y Obra).
\NewDocumentEnvironment{cita}{ O{} O{} O{} +b }{%
  \begin{adjustwidth}{2cm}{0pt}
  \raggedright
  #4\par
  \raggedleft
  \if\relax\detokenize{#1#2#3}\relax
    % sin firma
  \else
    #1%
    \if\relax\detokenize{#2}\relax\else\if\relax\detokenize{#1}\relax\else\space\fi(#2)\fi
    \if\relax\detokenize{#3}\relax\else\if\relax\detokenize{#1#2}\relax\else, \fi\textit{#3}\fi
  \fi
  \end{adjustwidth}
}{}



%% ------------------- Listas----------------------%%
    % --- Lista alfabética ---
    \newenvironment{la}
      {\begin{list}{\alph{alphaitem})}{%
          \usecounter{alphaitem}%
          \setlength{\leftmargin}{1cm}%
          \setlength{\parskip}{3pt}% 
          \setlength{\itemsep}{0pt}% ← espacio entre ítems
          \setlength{\parsep}{0pt}%  ← párrafos dentro de un ítem
      }}
      {\end{list}}
      
    % --- Lista numérica ---
    \newenvironment{lnum}
      {\begin{list}{\arabic{numitem}.}{%
          \usecounter{numitem}%
          \setlength{\leftmargin}{1cm}%
          \setlength{\parskip}{3pt}% 
          \setlength{\itemsep}{0pt}% ← espacio entre ítems
          \setlength{\parsep}{0pt}%  ← párrafos dentro de un ítem
      }}
      {\end{list}}

%% ------------------- Preguntas Test ----------------------%%
      \usepackage{xparse} % si ya lo tienes, no lo repitas

    \NewDocumentEnvironment{test}{m m o}{%
      \begin{tcolorbox}[
        colback=white,
        colframe=black!15,
        boxrule=0.6pt,
        arc=2mm,
        left=2mm,right=2mm,top=2mm,bottom=2mm
      ]
      \centering
      \includegraphics[width=\linewidth]{#1}
      \end{tcolorbox}
    
      \vspace{2mm}
    
      % Caja SOLO para texto (SÍ partible y mantiene márgenes al partir)
      \begin{tcolorbox}[
        enhanced,
        breakable,
        colback=black!3,
        colframe=black!25,
        boxrule=0.6pt,
        arc=2mm,
        left=2mm,right=2mm,top=1.5mm,bottom=1.5mm
      ]
      \textbf{Respuesta:} \textbf{#2}%
      \IfValueT{#3}{\par\smallskip \textbf{Justificación:} #3}
      \end{tcolorbox}
    }{}
      
% ------------------- Imágenes----------------------%
    \graphicspath{{Img/}}% Rutas por defecto 
    % --- Imagen adaptativa ---
    % Registros de dimensión definidos una sola vez
    \newdimen\imgcap
    \newdimen\imgmin
    \newdimen\imgmax
    \newdimen\imgremain
    \newdimen\imgh
    \newdimen\imgneed
    
    \newcommand{\imagenfit}[3]{%
      \addvspace{10pt}% separación antes
      \par\begingroup
        % Parámetros de composición (asignaciones, no \newdimen)
        \imgcap=1\baselineskip            % alto estimado de título+fuente
        \imgmin=0.15\textheight
        \imgmax=0.15\textheight
        \imgremain=\pagegoal
        \ifdim\pagegoal=\maxdimen \imgremain=0pt\fi
        \advance\imgremain by -\pagetotal
        \advance\imgremain by -\imgcap
        % Decisión de altura destino
        \ifdim\imgremain<\imgmin
          \newpage
          \imgh=0.20\textheight
        \else
          \imgh=\imgremain
          \ifdim\imgh>\imgmax \imgh=\imgmax \fi
        \fi
        % Reservar espacio para evitar cortes del bloque completo
        \imgneed=\imgh \advance\imgneed by \imgcap
        \Needspace{\imgneed}
        % Cuerpo
        \noindent\begin{minipage}{\textwidth}
          \centering
          \captionof{figure}{\textemdash\ #2}\par\vspace{0.3em}% título arriba
          \includegraphics[width=\linewidth,height=\imgh,keepaspectratio]{\detokenize{#1}}\par
          \vspace{0.3em}\small\textit{Fuente: }#3% fuente abajo
        \end{minipage}%
      \endgroup
      \par\addvspace{10pt}%
    }

%------------ Bloque de RECUADRO ESPECIAL -------------%
    \newenvironment{specialblock}[1]{%
      \begin{quote} % crea sangría a izquierda y derecha
      \small % texto más pequeño
      \noindent\textbf{#1}\\[-0.3em] % título en negrita
      \rule{\linewidth}{0.4pt}\\[0.6em]}{
      \end{quote}}
      
%-------------------------- Portada -----------------------%
\newcommand{\oldtitle}[1]{\def\OldTitle{#1}}
\newcommand{\changerationale}[1]{\def\ChangeWhy{#1}}

\newcommand{\changebox}{%
\begin{tcolorbox}[title={Historial de cambios del tema}]
\textbf{Título anterior:} \OldTitle\par
\textbf{Motivo del cambio:} \ChangeWhy
\end{tcolorbox}
}

%-------------------------- Author Font -----------------------%
    \newcommand{\authorfont}[1]{{\fontfamily{EBGaramond-TLF}\selectfont #1}}

%-------------------------- Anexos -----------------------%
    \makeatletter
    \newcounter{anexo}
    \renewcommand{\theanexo}{\Roman{anexo}}
    \newcommand{\anexo}[1]{%
      \clearpage
      \phantomsection
      \refstepcounter{anexo}%
      \noindent\textbf{Anexo \theanexo.- }#1\par\medskip
      \addcontentsline{stc}{section}{Anexo \theanexo.- #1}%
    }
    \makeatother

    %-------------------------- Lista de figuras (personal) -----------------------%
\makeatletter
\newcommand{\MyListOfFigures}{%
  \clearpage
  \phantomsection
  {\centering\bfseries\Large Índice de figuras\par}%
  \medskip
  % Entrada en nuestro índice de contenido (stc)
  \addcontentsline{stc}{section}{Índice de figuras}%
  % Recuperación del listado (sin cabecera propia)
  \begingroup
    \parindent=0pt\parskip=2pt
    \@starttoc{lof}%
  \endgroup
}
\makeatother

%-------------------------- Bibliografía -----------------------%
\makeatletter
% Contador para las referencias
\newcounter{myref}
% Cabecera de la bibliografía, entra en el índice 'stc'
\newcommand{\MyBibliography}{%
  \clearpage
  \phantomsection
  {\centering\bfseries\Large Bibliografía y referencias\par}%
  \medskip
  \addcontentsline{stc}{section}{Bibliografía y referencias}%
  \setcounter{myref}{0}%
}
\newcommand{\MyBibItem}[4]{%
  \refstepcounter{myref}%
  \noindent[\themyref]\ #1.\ \emph{#2}. #3. #4\par
}
\makeatother

%--------- Establecer cuatro niveles de índice --------%
    \setcounter{tocdepth}{4}   
    \setcounter{secnumdepth}{4}% numera hasta \paragraph

%%%%%%%%%%%%%%%%%%%%%%%%%%%%%%%%

\makeatletter

    %--------- Interlineado Global --------%
    \edef\GlobalBaseLineStretch{\baselinestretch}
    
    %--------- Espaciado de seccinones --------%
    \renewcommand\section{\@startsection{section}
      {1}{\z@}
      {2.5ex \@plus 1ex \@minus .2ex} % espacio antes
      {1.5ex \@plus .2ex}             % espacio después
      {\normalfont\Large\bfseries}    % formato del título
    }
    \renewcommand\subsection{\@startsection{subsection}{2}{\z@}
      {3ex \@plus 0.8ex \@minus .2ex}
      {1.5ex \@plus .2ex}
      {\normalfont\large\bfseries}
    }
    \renewcommand\subsubsection{\@startsection{subsubsection}{3}{\z@}
      {3ex \@plus 0.5ex \@minus .2ex}
      {1ex \@plus .2ex}
      {\normalfont\normalsize\bfseries}
    }
    \renewcommand\paragraph{\@startsection{paragraph}{4}{\z@}%
    {-2ex \@plus -0.5ex \@minus -.2ex}% espacio antes
    {0.8ex \@plus .2ex}% espacio después
    {\normalfont\normalsize\itshape}}% ← cursiva en lugar de negrita

    %--------- Ecuaciones --------%   
    % Variables globales para almacenar títulos y notas
    \gdef\leq@current@section{}%
    \gdef\leq@current@subsection{}%
    \gdef\leq@last@printed@section{}%
    \gdef\leq@last@printed@subsection{}%
    
    % Comando para añadir nota (invisible en el cuerpo)
    \newcommand{\eqnote}[1]{%
      \addcontentsline{leq}{leqnote}{#1}%
    }
    
    % Comando para ecuaciones
    \newcommand{\eqblock}[2]{%
      % Verificar si necesitamos imprimir el título de sección
      \ifx\leq@current@section\leq@last@printed@section\else
        \addcontentsline{leq}{leqsection}{\leq@current@section}%
        \global\let\leq@last@printed@section\leq@current@section
        \gdef\leq@last@printed@subsection{}% Reset subsección al cambiar sección
      \fi
      % Verificar si necesitamos imprimir el título de subsección
      \ifx\leq@current@subsection\leq@last@printed@subsection\else
        \ifx\leq@current@subsection\@empty\else
          \addcontentsline{leq}{leqsubsection}{\leq@current@subsection}%
          \global\let\leq@last@printed@subsection\leq@current@subsection
        \fi
      \fi
      % Ecuación
      \refstepcounter{eqnum}%
      \begin{equation*}
        #1\tag{E.\theeqnum}
      \end{equation*}%
      \addcontentsline{leq}{leqt}{[E.\theeqnum] -- #2}%
      \addcontentsline{leq}{leqm}{#1}%
    }
    
    %%% ----- Índice de ecuaciones ----------%%%
    % Tamaño para []
    \newcommand{\leqIndexFont}{\normalsize}
    \newcommand{\leqTitleFont}{\tiny}
    \newcommand{\leqNoteFont}{\tiny}
    % Cabecera de sección 
    \newcommand*\l@leqsection[2]{%
      \addvspace{25pt}% Mayor separación antes de nueva sección
      \noindent\leqIndexFont
      \makebox[\linewidth]{\rule{.38\linewidth}{0.4pt}\ \ \textbf{  \MakeUppercase{#1}}\ \ \rule{.38\linewidth}{0.4pt}}%
      \par\vspace{-10pt}
    }
    % Cabecera de subsección 
    \newcommand*\l@leqsubsection[2]{%
      \par\addvspace{15pt}%
      \noindent\leqIndexFont #1
      \par\vspace{-7pt}
    }
    % Título de ecuación 
    \newcommand*\l@leqt[2]{%
    \par\addvspace{10pt}%
      \par\noindent\leqTitleFont #1\par
    }
    % Miniatura de ecuación
    \newcommand*\l@leqm[2]{%
      \vspace{0.3\baselineskip}%
      \par\scriptsize\noindent\hspace*{1.5em}\(#1\)\par%
    }
    % Nota de ecuación 
    \newcommand*\l@leqnote[2]{%
      \vspace{0.2\baselineskip}%
      \par\noindent\leqNoteFont\hspace*{1.5em}\textbf{Nota.--} #1\par\vspace{0.5\baselineskip}%
    }
    % Generación del índice
    \newcommand{\mylistofequations}{%
      \clearpage
      \phantomsection
      % Título visual de la lista (ajústalo a tu gusto):
      {\centering\bfseries\Large Índice de ecuaciones\par}
      % Entrada en nuestro índice 'stc':
      \addcontentsline{stc}{section}{Índice de ecuaciones}%
      % Llamada a tu lista real (usa el nombre exacto de tu macro):
      \@starttoc{leq}%
    }
    
    %--------- Captura de títulos de sección --------%
    \let\leq@original@section\section
    \let\leq@original@subsection\subsection
    % Flag para saber si estamos en una sección asterisco
    \newif\ifLeqStarSection
    \LeqStarSectionfalse
    
    % Redefinir \section CON CONTROL DE ASTERISCO
    \renewcommand{\section}{%
      \@ifstar{\leq@section@star}{\@dblarg\leq@section@nostar}%
    }
    \def\leq@section@star#1{%
      \global\LeqStarSectiontrue
      \gdef\leq@current@section{#1}%
      \gdef\leq@current@subsection{}%
      \leq@original@section*{#1}%
      \global\LeqStarSectionfalse
    }
    \def\leq@section@nostar[#1]#2{%
      \global\LeqStarSectionfalse
      \gdef\leq@current@section{#2}%
      \gdef\leq@current@subsection{}%
      \leq@original@section[#1]{#2}%
    }
    % Redefinir \subsection
    \renewcommand{\subsection}{%
      \@ifstar{\leq@subsection@star}{\@dblarg\leq@subsection@nostar}%
    }
    \def\leq@subsection@star#1{%
      \global\LeqStarSectiontrue
      \gdef\leq@current@subsection{#1}%
      \leq@original@subsection*{#1}%
      \global\LeqStarSectionfalse
    }
    \def\leq@subsection@nostar[#1]#2{%
      \global\LeqStarSectionfalse
      \gdef\leq@current@subsection{#2}%
      \leq@original@subsection[#1]{#2}%
    }

    % ----- Restauración de valores Globales de estilo ----------%
    % Flags
    \newif\ifInFootnote
    \InFootnotefalse
    \pretocmd{\@footnotetext}{\InFootnotetrue}{}{}
    \apptocmd{\@footnotetext}{\InFootnotefalse}{}{}
    % Profundidad de listas (soporta anidadas)
    \newcount\MyListDepth \MyListDepth=0
    \AtBeginEnvironment{la}{\advance\MyListDepth by 1}
    \AfterEndEnvironment{la}{\advance\MyListDepth by -1}
    \AtBeginEnvironment{lnum}{\advance\MyListDepth by 1}
    \AfterEndEnvironment{lnum}{\advance\MyListDepth by -1}
    \AtBeginEnvironment{itemize}{\advance\MyListDepth by 1}
    \AfterEndEnvironment{itemize}{\advance\MyListDepth by -1}
    \AtBeginEnvironment{enumerate}{\advance\MyListDepth by 1}
    \AfterEndEnvironment{enumerate}{\advance\MyListDepth by -1}
    % Restauración diferida: hazlo al INICIO del próximo párrafo real
    \newif\if@pendingRestore
    \@pendingRestorefalse
    \newcommand{\RequestRestore}{\global\@pendingRestoretrue}
    \everypar{%
      \if@pendingRestore
        \global\@pendingRestorefalse
        \ifInFootnote\else
          \ifnum\MyListDepth=0
            \setlength{\parskip}{\GlobalParskip}%
            \setstretch{\GlobalBaseLineStretch}%
          \fi
        \fi
      \fi
    }
    % Al final de bloques:
    \AfterEndEnvironment{equation}{\RequestRestore}
    \AfterEndEnvironment{equation*}{\RequestRestore}
    \AfterEndEnvironment{la}{\RequestRestore}
    \AfterEndEnvironment{lnum}{\RequestRestore}
    % Al final de macros:
    \apptocmd{\eqblock}{\RequestRestore}{}{}
    \apptocmd{\imagenfit}{\RequestRestore}{}{}
    
\makeatother

% ===================== ToC propio con 4 niveles (único bloque) =====================
\makeatletter

% --- Escritores: añadimos una línea al .stc en cada cabecera numerada ---
% Guardamos las versiones actuales para llamar al comportamiento normal
\let\MyTOC@orig@section\section
\let\MyTOC@orig@subsection\subsection
\let\MyTOC@orig@subsubsection\subsubsection
\let\MyTOC@orig@paragraph\paragraph

% SECTION
\renewcommand{\section}{%
  \@ifstar{\MyTOC@section@star}{\@dblarg\MyTOC@section@nostar}%
}
\def\MyTOC@section@star#1{%
  \MyTOC@orig@section*{#1}% sin entrada al .stc
}
\def\MyTOC@section@nostar[#1]#2{%
  \MyTOC@orig@section[{#1}]{#2}% comportamiento normal (escribe al .toc)
  % Escribimos también nuestra línea al .stc (texto con \numberline)
  \addcontentsline{stc}{section}{\protect\numberline{\thesection}#2}%
}

% SUBSECTION
\renewcommand{\subsection}{%
  \@ifstar{\MyTOC@subsection@star}{\@dblarg\MyTOC@subsection@nostar}%
}
\def\MyTOC@subsection@star#1{%
  \MyTOC@orig@subsection*{#1}%
}
\def\MyTOC@subsection@nostar[#1]#2{%
  \MyTOC@orig@subsection[{#1}]{#2}%
  \addcontentsline{stc}{subsection}{\protect\numberline{\thesubsection}#2}%
}

% SUBSUBSECTION
\renewcommand{\subsubsection}{%
  \@ifstar{\MyTOC@subsubsection@star}{\@dblarg\MyTOC@subsubsection@nostar}%
}
\def\MyTOC@subsubsection@star#1{%
  \MyTOC@orig@subsubsection*{#1}%
}
\def\MyTOC@subsubsection@nostar[#1]#2{%
  \MyTOC@orig@subsubsection[{#1}]{#2}%
  \addcontentsline{stc}{subsubsection}{\protect\numberline{\thesubsubsection}#2}%
}

% PARAGRAPH
\renewcommand{\paragraph}{%
  \@ifstar{\MyTOC@paragraph@star}{\@dblarg\MyTOC@paragraph@nostar}%
}
\def\MyTOC@paragraph@star#1{%
  \MyTOC@orig@paragraph*{#1}%
}
\def\MyTOC@paragraph@nostar[#1]#2{%
  \MyTOC@orig@paragraph[{#1}]{#2}%
  % Si no numeras los \paragraph, cambia \theparagraph por texto plano sin \numberline
  \addcontentsline{stc}{paragraph}{\protect\numberline{\theparagraph}#2}%
}

% --- Impresor del índice propio (lee el .stc) ---
\newcommand{\MyTOC}{%
  \par\bigskip
  {\centering\bfseries\Large Índice de contenido\par}%
  \medskip
  \begingroup
    \parindent=0pt\parskip=2pt
    \@starttoc{stc}%
  \endgroup
}

\makeatother
% ===================== fin ToC propio =====================


%%%%%%%%%%%%%%


% --- Datos del tema ---
\title{Tema 3.A.34 -- Teorías de la inversión en bienes de equipo.\\
\large Incertidumbre e irreversibilidad de la inversión. Implicaciones de política económica}
\author{Marcelo Ortega}
\date{Marzo 2026}
\newcommand{\TemaCode}{3A34}

\oldtitle{Tema 3.A.33. Teorías de la demanda de inversión. Implicaciones de política económica}
\changerationale{
    Se provee de más detalle acerca de lo que se espera que el opositor desarrolle en su exposición.
}

\begin{document}


\maketitle
\changebox

\newpage

% --- Página de resumen y modificaciones ---
\puntosclave{ %% Escribir puntos clave Apuntes CECO

    }
\vspace{1cm}

\modificaciones{
    
    }

\newpage
\MyTOC
\newpage

\mylistofequations
\clearpage

\MyListOfFigures
\clearpage


\pagenumbering{arabic}
\setcounter{page}{1}


\section{Introducción}\label{intro}


\subsection{Contextualización}\label{context}


\subsubsection{Relevancia}\label{importance}


\subsubsection{Historia}\label{history}



\subsubsection{Problemática}\label{problem} 




\subsection{Estructura de la exposición}\label{structure}








\clearpage
\section{Teoría de la Inversión: Información perfecta}
\puntosclave{
    A partir de la inclusión de costes de ajustes en el modelo de Jorgenson, se observa que $q_t$ es la variable que sintetiza toda la información relevante sobre el futuro que la empresa necesita conocer para adoptar sus decisiones de inversión.
    
    La teoría de la $q$ de Tobin estudia la relación entre el valor de mercado y el coste de reposición del capital. 
    
    La construcción del diagrama de fases pone de manifiesto que los efectos de política económica serán diferentes dependiendo de si el cambio es permanente o temporal.
}


A lo largo de esta primera sección vamos a tratar de exponer los principales componentes/características del marco teórico dominante acerca de la Inversión en bienes de equipo. Debemos tener presente que, cualquiera que sea el tipo de modelo que se utilice, para determinar la demanda de inversión en capital se ha de prestar atención a tres elementos:
\begin{lnum}
    \item El volumen de capital productivo que la empresa juzga deseable tener en un momento dado del tiempo;
    \item El ritmo o secuencia temporal al que la empresa proyecta cubrir la diferencia entre el volumen de capital deseado y el existente (lo que determina la inversión neta); y,
    \item La secuencia de inversión de reposición (lo que determina la inversión bruta).
\end{lnum} 

Estos tres elementos se desprenden de un razonamiento muy sencillo: \textbf{las empresas comprometen sus posibilidades de producción futura a sus decisiones de inversión actual}. De esta forma, el nivel de inversión óptimo debe responder a tres cuestiones esenciales: \textbf{¿cuánto?, ¿cuándo? Y, ¿en qué?}.\footnote{%
    Cuánto y en qué son cuestiones estrechamente relacionadas con el destino de la inversión: la empresa puede decidir mantener su flujo de inversión constante, lo que no ampliará su stock de capital futuro pero impedirá la desamortización de capital derivada de la depreciación experimentada. Y, por otro lado, cuándo está estrechamente relacionado con las condiciones frente a las cuáles las empresas consideran óptimo cambiar su política de inversión en bienes de equipo.
} Una versión primitiva que permite dar respuesta a estas preguntas la podemos sintetizar perfectamente en la siguiente formulación análitica:

\eqblock{
\begin{aligned}
    \boxed{\begin{aligned}
        &\max_{I_t^o}\quad {VAN}_t=\int_0^\infty e^{-rt}[P_xF(K_t)-rI_t^o]\mathrm dt\\[0.5em]
        &\text{s.a.}\quad \underset{\text{\scriptsize Inversión neta = (Bruta + depreciación)}}{I_t=\dot K_t+\delta K_t}
    \end{aligned}}\Longrightarrow\not\exists\text{Optimización dinámica}
\end{aligned}
}{Problema de optimización del productor: Inversión neta. Teorías de la Inversión en bienes de equipo: certidumbre}

Este sencillo problema de optimización de la empresa representativa consiste por tanto en la maximización del valor presente de sus beneficios (o, más bien, de los rendimientos de los proyectos de inversión) sujeto a una restricción de acumulación de capital:

inversión neta = inversión bruta + depreciación del capital

\textcolor{magenta}{(¿qué quiere decir esto? No entiendo muy bien) Ahora bien, las respuestas ante un cambio favorable, por ejemplo,  en el precio del capital serían: (i) tanto como para volver a la condición marginal óptima (rentabilidad-coste, implica oferta elástica); (ii) ya, ¿por qué esperar?; y, (iii) en mantener y aumentar nuestro stock de capital instalado.}


Esto es consecuencia directa de una limitación esencial: este problema no tiene ningún rasgo dinámico per se. La empresa bien podría ejecutar este mismo problema de optimización cada periodo y alcanzar las mismas conclusiones de una manera más sencilla ya que no existe diferencia entre planificación y ejecución. Pero, ¿no piensan (deberian pensar) los directores en el futuro? ¿es que no llevan a cabo una labor planificadora? Más aún, ¿no existen restricciones en el mercado de capital? ¿es la oferta de capital perfectamente elástica?

Evidentemente, todas estas respuestas conducen a exigir la introducción de un elemento que otorgue el carácter intertemporal propio que acarrean las decisiones de inversión.

Existen diversas maneras de introducir este elemento intertemporal a las decisiones de inversión de los productores. Una de las aproximaciones pioneras fue la de JORGENSON (1963), bajo el marco de la SNC. 

Si bien es cierto que JORGENSEN lleva a cabo una adaptación \textit{ad hoc}, es el primero en enfatizar la importancia de los retardos y restricciones que experimentan las decisiones de inversión.\footnote{%
    Para un repaso histórico de la Teoría de la Inversión de bienes de equipo [\authorfont{Ver Anexo \ref{historia_teoria}}]: KEYNES, SNC.
} 

\subsection{Costes de Ajuste: ¿de quién es?}
La falta de microfundamentación de dicha especificación, basada en reatrdos parámetricos construidos ad hoc, exigía su reformulación para inmunizarse frente a la revolución metodológica posterior a la Crítica de LUCAS.

Una opción es considerar la existencia de costes de ajuste del stock de capital. 

\textcolor{magenta}{¿De quién fue la teoría de costes de ajuste? (formulación que se propone aquí)}


\subsubsection{Tipología: ¿qué tipos de costes de ajuste hay?}
Estos costes pueden ser de dos tipos: internos o externos.

\textcolor{magenta}{continuar (desarollar algo aquí que se ainteresante/relevante sobre los costes y que no tenga estrecha relación ni con los internos ni con los externos).... Alargar esta introducción un poco, en definitiva.}

\paragraph{Costes de ajuste: Internos}
Los costes de ajuste internos son los costes directos en los que incurren las empresas cuando modifican su stock de capital. Por ejemplo, la instalación de nuevo capital o la formación de los trabajadores en el manejo de nueva maquinaria. 

\textbf{¿INTRODUCIR GRÁFICO DE INVERSIÓN - COSTES DE AJUSTE INTERNOS?}

Para verlo de forma intuitiva, podemos dibujar el plano con el eje \textcolor{magenta}{...} y el eje \textcolor{magenta}{..}. Imaginemos que sucedería ante \textcolor{magenta}{... (continuar aquí con un ejemplo quese puede visualizar osbre el gráfico.) Este era el anterior: Sabemos que, a medida que la tasa de variación del stock de capital se acerca a infinito ($\dot K_t\to \infty$), los costes de ajuste también lo hacen ($C(I_t)\to\infty$). Por tanto, una disminución del tipo de interés llevaría a un incremento de la demanda de inversión por una mejora en los precios/costes relativos ($\downarrow r_t\to\; I_t^o\uparrow$) de los bienes de equipo.}

En definitiva, los costes de ajuste internos nos permiten ver cómo se ajusta la variación dinámica del capital ante sucesos favorables/desfavorables: \textbf{el stock de capital aumentará/descenderá paulatinamente hacia el nuevo nivel deseado}. 

\paragraph{Costes de ajuste: Externos}
Los costes de ajuste externos, por su parte, surgen como consecuencia de mercados de capitales imperfectos. Por ejemplo, si la empresa quiere disponer de capital con rapidez, el coste de financiación será mayor. 

\textbf{¿INTRODUCIR GRÁFICO DE INVERSIÓN - COSTES DE AJUSTE INTERNOS?}

Para verlo de forma intuitiva, podemos dibujar el plano con el eje \textcolor{magenta}{...} y el eje \textcolor{magenta}{..}. Imaginemos que sucedería ante \textcolor{magenta}{... (continuar aquí con un ejemplo quese puede visualizar osbre el gráfico.) Este era el anterior: Intuititavemnte, cuando la oferta de capital no es perfectamente elástica, una variación discreta que dé lugar a un aumento del stock deseado de capital eleva el precio de los bienes de capital.... (incompleto, no se observa bien).} 

\subsubsection{Desarrollo: ¿...?}
Incorporando ambos tipos de funciones al programa de optimización del productor es una tarea relativamente sencilla. \textcolor{magenta}{(footnote: Se relega a los anexos el desarrollo analítico)}

Para facilitar el análisis de sus implicaciones de política económica y el posterior análisis en contexto de incertidumbre impondremos dos supuestos. En primer lugar, que la función de producción exhibe rendimientos constantes de escala, los mercados de productos son competitivos y la oferta de factores, a excepción de la de capital, es perfectamente elástica. Esto permite que los beneficios reales de una empresa representativa en el período $t$ sean proporcionales a su stock de capital. 

En segundo lugar, asumimos una función de costes convexa. Esto implica que se penalizan cambios rápidos en las variables agregadas, favoreciéndose movimientos graduales en el tiempo, de forma que el coste marginal de ajuste aumenta cuanto mayor es la variación del stock de capital, y adquieriendo el problema una dinámica temporal muy relevante: \textbf{las empresas siempre incurren en menores costes de ajuste si distribuyen un proyecto de inversión a lo largo de más años}. Como vemos, este supuesto tiene fuertes implicaciones para el comportamiento de la inversión y no es totalmente consistente con el comportamiento observado. Si este fuera el caso, entonces esperaríamos ver una inversión continua y gradual por parte de las empresas en lugar de grandes y discretas adiciones de capital.\footnote{%
  Una serie de influyentes estudios de equilibrio parcial (CABALLERO y ENGEL 1999; COOPER et al. 1999; CABALLERO et al. 1995) han argumentado que los modelos de inversión con costes no convexos superan empíricamente superan a los modelos de costes convexos porque pueden de costes convexos, ya que pueden producir cambios desproporcionadamente bruscos en la demanda de demanda de inversión después de grandes shocks a nivel agregado.
}

\textcolor{magenta}{(En el texto original, se asumía esta tercer condición. No entiendo muy bien porqué, por lo que, idealmente deberíamos poder descartarla siempre y cuando no genera conflicto con el desarrollo posterior). \textbf{$\not\exists$ Costes externos}. En primer lugar, supondremos que estos costes de ajuste son únicamente internos (es decir, el precio de compra de los bienes de capital es constante e igual a 1), aunque es fácil reinterpretar el modelo suponiendo costes de naturaleza externa (ABEL (1982), HAYASHI (1982) y SUMMERS (1981b))}

\textcolor{magenta}{((Explicar aquí con palabras el planteamiento del problema, su desarrollo y resolución)) Esto es lo que estaba, expandir: Para la resolución del problema en tiempo continuo planteamos simplemente el HAMILTONIANO del valor presente, introduciendo la restricción en el problema de maximización, y sustituyendo el multiplicador dinámico ($\lambda$) en términos presentes a partir de la siguiente expresión $q_t=\lambda e^{-rt}$, obtenemos:}

Las condiciones que obtendremos serán:

\eqblock{
  \begin{aligned}
    &\underset{\text{\scriptsize \textbf{invertir} hasta condición marginal}}{\boxed{q_t=r+C'(I_t^o)}}\\[1em]
    &\underset{\text{\scriptsize \textbf{acumular} hasta condición marginal}}{\boxed{P_xF'(K_t)=(r+\delta)q_t-\dot q_t}}\\[1em]
    &\underset{\text{\scriptsize \textbf{dinámica} del capital}}{\boxed{I_t^o-\dot K_t-\delta K_t=0}}
  \end{aligned}
}{}

Estas condiciones pueden fácilmente si se les proporciona una intución económica. En primer lugar, la \textbf{condición de óptimo en la inversión} nos permite conocer hasta cuando invertirá la empresa: hasta que el valor adicional esperado que genere una unidad adicional de capital se iguale al coste de adquisición de dicha unidad de capital. \textcolor{magenta}{Explicar más: ¿por qué? ¿qué significa, etc?}

En segundo lugar, la \textbf{condición de óptimo en la acumulación} nos permite conocer dónde se encuentra el nivel óptimo de capital para la empresa: donde el ingreso marginal que genere una unidad adicional de capital en la empresa se iguale al coste de oportunidad de una unidad de capital.\textcolor{magenta}{Explicar más: ¿por qué? ¿qué significa, etc?}

En tercer, y último, lugar, la \textbf{condición dinámica} permite recuperar la restricción del problema y garantizar su cumplimiento estricto.\textcolor{magenta}{Explicar más: ¿qué quiere decir? ¿por qué? ¿qué significa, etc?}
  
\textcolor{magenta}{(¿Esto es una CPO o es un supuesto previo? La misma pregunta esta en el Anexo: pensar, resolver, y situar donde corresponda. )Por último, la \textbf{condición de transversalidad} evita que se generen sendas explosivas de inversión y, por tanto, dotar de sentido económico al problema. Así, esta condición establece que el límite del producto de la variable coestado descontada y la variable estado debe ser igual a cero, rechazándose que un comportamiento puramente especulativo sea óptimo para la empresa (en tanto la única razón para mantener el activo correspondiente al valor marginal del capital es la expectativa de ganancias de capital superiores e indefinidas, no ligadas a los beneficios actualmente derivados de su uso en la producción)}

Como podemos ver, la variable $q_t$, \textcolor{magenta}{que es.... (explicar qué es y de dónde sale, cómo se obtiene)} . Por tanto, sintetiza toda la información relevante sobre el futuro que la empresa necesita conocer para adoptar sus decisiones de inversión. ¿Por qué? Poruqe $q$ refleja los efectos que tendría una unidad adicional de capital sobre el valor presente de los beneficios futuros: es el valor presente descontado de los ingresos marginales futuros de una unidad adicional de capital. 

\subsubsection{Estabilidad del equilibrio: Unicidad de la senda de convergencia}
Recuperando dos de las ecuaciones expuestas, representar la dinámica de la inversiónmediante un diagrama de fases que nos permite incidir en el análisis de la interacción de la inversión con los instrumentos de Política Económica:

\eqblock{
\begin{aligned}
    &\begin{aligned}
      &\overset{\dot K_t|_{K^*}=0}{\underset{\text{\scriptsize equilibrio}}{\Longrightarrow}}\quad {q_{EE}^*=1}\\[0.5em]
      &\underset{\delta=0}{\Longrightarrow}\quad\dot q_t=rq_t-P_xF'(K_t)
    \end{aligned}\overset{\scriptsize{\text{equilibrio}}}{\Longrightarrow}
    \left\{\begin{aligned}
        &(1)&&\hspace{3em}\underset{\text{\scriptsize el nivel de inversión es óptimo}}{\boxed{q_{EE}^*=1}}\\[1em]
        &(2)&&\underset{\substack{\text{\scriptsize dinámica del valor sombra del capital adicional}}}{\boxed{\dot q_t=0\iff rq_t=P_xF'(K_t)}}
    \end{aligned}\right.
\end{aligned}
}{Condiciones dinámicas óptimas: Diagrama de fases. Teoría de la Inversión en bienes de equipo: Certidumbre}

Desde esta perspectiva, parece claro que la convergencia hacia al equilibrio dependerá del valor que adquiera la $q_t$: \textbf{siempre que el valor descontado al presente de los rendimientos de una inversión sean mayores a los costes de adquirir dicho capital nuevo (en el margen), la inversión debe realizarse}. 

Estas ecuaciones y condiciones de equilibrio nos permiten, como hemos anticipado, abordar el análisis gráfico del modelo a partir de la construcción de un diagrama de fases en un espacio $(K, q)$ y mostrar las sendas óptimas de inversión en el tiempo:

\imagenfit{ICEX_Fases_Q.png}{(Dcha.) Comportamiento del stock de capital | (Izq.) Comportamiento de $q$}{\textit{Macroeconomía avanzada}. ROMER (2006)}

Como vemos, la \textbf{primera condición} \textcolor{magenta}{viene representada por...} y muestra el comportamiento del stock de capital a partir de la primera CPO.\footnote{%
  De acuerdo con la $q$ de TOBIN, que luego veremos, en el equilibrio el mercado de valores estaría concediendo al capital instalado en la empresa un valor igual a su coste de reposición, de forma que la empresa no tendría incentivos a invertir o desinvertir. 
} Por otro lado, la \textbf{segunda condición} \textcolor{magenta}{viene representada por ...} y muestra el comportamiento de la $q$ a partir de la segunda CPO.

Combinando los dos gráficos, obtenemos la división del plano en cuatro regiones donde conocemos perfectamente la dinámica de ambas variables estado:

\imagenfit{ICEX_Fases_Conjunto.png}{(Dcha.) Diagrama de fases | (Izq.) Sendero hacia el punto de inflexión (saddle-point)}{\textit{Macroeconomía avanzada}. ROMER, 2006}

Gráficamente: (i) dos regiones divergentes en las que no se cumple la condición de transversalidad y, por tanto, conducen a la empresa hacia sendas de inversión explosivas o implosivas no tendentes al estado estacionario (punto E); y, (ii) dos regiones convergentes donde sí se cumple la condición de transversalidad y determinan la trayectoria que conduce al estado estacionario de forma estable. 

Por lo tanto, esas regiones no contienen sendas estables sino que contienen posibles trayectorias del sistema. \textbf{De todas esas trayectorias posibles, solo una (la que coincide exactamente con el sendero de silla) converge al equilibrio.} Y este es precisamente la clave del modelo: dado un nivel de capital instalado $K_0$, existe una única senda estable que conduce al equilibrio, y en consecuencia un único $q_0$ que nos sitúa dentro de esta trayectoria. 

\textcolor{magenta}{(explicar mejor por qué esto sucede así, he eliminado la ecuación de formulación analítica, pasándola a los Anexos): (Antes esta era la explicación) La razón profunda es estructural. Por un lado, $K$ es una variable predeterminada que no puede adaptarse inmediatamente/saltas de un punto a otro. Por otro lado, $q$ es una variable de salto, que permite adaptarse al sistema a las nuevas condiciones estructurales.}

Por tanto, la lógica del problema será: dado el capital heredado con el que arranca la empresa, el mercado debe asignar inmediatamente a la empresa exactamente el valor de $q$ que la sitúa sobre la senda estable. Si le asigna un valor demasiado bajo, aunque sea todavía superior a uno, la inversión inicial será positiva pero insuficiente en relación con la trayectoria óptima, y la dinámica posterior de $q$ impedirá la convergencia. Por el contrario, si le asigna un valor demasiado alto la inversión será excesiva y la trayectoria también se alejará del equilibrio. En ambos casos se acaba violando la condición de transversalidad.

\textcolor{magenta}{(revisar, creo que estamos diciendo la misma cosa dos veces, poner como nota en rojo la correción y la explicación de por qué si/no se está diciendo correctamente)Por tanto, en equilibrio se cumplen dos condiciones. En primer lugar, que el \textbf{capital no cambia: la inversión neta es nula}. Como ya se ha señalado anteriormente, que $q$ sea igual a $1$ significa que el valor descontado de los beneficios futuros por un euro adicional de capital instalado sea igual al valor de adquisición de una unidad de capital: las empresas no tienen ningún incentivo para aumentar o reducir su stock de capital. En segundo lugar, que el \textbf{valor sombre del capital no cambia: la condición de marginalidad es óptima}. Asimismo, para que $\dot q_t$ sea igual a $0$ cuando $q$ es igual a $1$ el ingreso marginal del capital debe ser igual a $r$, lo que implica que los beneficios asociados a la posesión de una unidad de capital compensan el valor de los intereses que se dejan de percibir (los inversores están satisfechos de mantener un capital del que no esperan obtener ni ganancias ni pérdidas).}

Son dos caras de la misma moneda: la situación es \textbf{óptima en capital instalado}. \textcolor{magenta}{(Reformular atendiendo a los cambios realizados: la Q de TOBIN se plantea después) Por tanto, la ventaja de la $q$ de TOBIN como medida del incentivo para invertir reside en que refleja la rentabilidad futura esperada del capital así, como la actual, es decir, se trata de un enfoque forward-looking.}

\subsection{Implicaciones de Política Económica}
Por tanto, traslandolo al ámbito de la Política Económica, la Teoría de la inversión hace hincapié en que las decisiones de inversión dependen no sólo de la Política Económica actual sino también de la que los agentes esperan que se adopte en el futuro. Por ejemplo, el anuncio de una rebaja fiscal sobre los beneficios empresariales a partir del año siguiente conllevará un aumento esperado de los beneficios para la empresa. Este impacto positivo futuro tendrá un efecto alcista sobre el valor de mercado de las acciones de la empresa, aumentando el valor de la q de Tobin y, por tanto, fomentando la inversión en el presente.

Una vez armados con el instrumental gráfico: ¿qué podemos decir sobre la reacción de los agentes ante cambios en la Política Económica? Estudiemos varios casos. 

\subsubsection{Políticas Fiscal: Demanda Efectiva ($y_t$)}
Veamos, en primer lugar, qué sucedería ante un aumento de la demanda agregada, que incrementa los beneficios para un stock de capital dado y, progresivamente, eleva la producción de la industria.

\paragraph{Perturbación permanente: Infrainversión}
Si este perturbación fuese permanente, partiendo de una situación inicial de equilibrio en el largo plazo (punto $E$), se produciría un desplazamiento hacia arriba de la curva $\dot q=0$: para el mismo stock de capital, los beneficios son ahora más elevados, los inversores necesitan unas ganancias de capital menores para justificar su tenencia de activos empresariales. 

\imagenfit{ShockProduccion_TempYPerm.png}{Efectos de un aumento temporal y permanente de la producción agregada ($\uparrow y_t$)}{\textit{Macroeconomía avanzada}. ROMER, 2006}

Este desplazamiento de la curva $\dot q=0$, como sabemos, implica que $q$ salta inmediatamente hasta un punto situado sobre el nuevo sendero de silla existente para el stock de capital dado. 

A continuación, $K$ y $q$ se desplazan hacia abajo a lo largo de esta senda hacia el nuevo punto de equilibrio (punto $E'$). Como la tasa de variación del stock de capital es una función creciente de $q$, esto significa que $K$ aumenta súbitamente en el momento de producirse el cambio y regresa luego gradualmente a cero.

En conclusión, \textbf{un aumento permanente de la producción provoca que la inversión se eleve temporalmente y que se alcance un equilibrio final con un mayor stock de capital agregado}. 

La lógica es simple: cuando la demanda del bien que produce la industria se eleva, la producción aumenta y, como el stock de capital no puede ajustarse automáticamente, la rentabilidad del capital existente incrementa. A su vez, esta mayor rentabilidad atrae nuevas inversiones, de modo que el stock de capital comienza a aumentar. A medida que lo hace, la producción de la industria aumenta y el precio del producto comienza a disminuir, de modo que tanto los beneficios como el valor del capital caen. El proceso continúa hasta que el valor de un incremento marginal del capital se iguala a su coste de adquisición, momento a partir del cual el incentivo que provocaba nuevas inversiones ha desaparecido. 

\paragraph{Perturbación transitoria: Suavización de la senda de convergencia ($\dot q_t=0$)}
Si, por el contrario, el aumento fuese temporal, no puede haber un salto anticipado de $q$: si se produce un salto anticipado de $q$ hacia abajo, los propietarios de acciones empresariales experimentarán en ese momento, con certeza, pérdidas de capital a una tasa infinita, lo que significa que no habrá nadie dispuesto a conservarlas. 

Por el contrario, ante un cambio transitorio de la demanda, la empresa se situará en un punto $A$, superior al equilibrio inicial (punto $E$), pero no llega a situarse en el nuevo sendero de silla como ocurría ante un incremento permanente. Progresivamente, $q$ y $K$ se desplazan gradualmente hacia el punto $B$, a partir del cual las empresas inician un proceso de desinversión.\footnote{%
    Cabe destacar que la senda que recorren $K$ y $q$ cruza la línea $\dot K_t=0$ antes de situarse sobre el sendero de silla original, es decir, el stock de capital comienza a reducirse antes de que la producción retorne a su nivel inicial. Esto se debe a que el ajuste de capital genera costes y al quedar poco tiempo para que los beneficios se mantengan altos, existen ventajas y casi ningún inconveniente en que la reducción comience inmediatamente.
}

En conclusión, un aumento transitorio de la demanda agregada provoca un aumento de la inversión menor que en el caso de un aumento permanente, en tanto $q$ aumenta en menor medida (recordemos que el comportamiento de la inversión depende de $q$). No obstante, a diferencia del aumento permanente, no se alcanza un nuevo equilibrio a largo plazo con un mayor stock de capital agregado. 

Estos resultados ponen de manifiesto que los efectos sobre la inversión no solo se explican a partir del nivel actual de la producción sino de su trayectoria en el tiempo, en tanto la inversión es mayor cuando la producción acaba de aumentar que cuando lleva siendo alta durante un período prolongado.

\subsubsection{Política Monetaria: Tipos de interés}
Los efectos de una Política Monetaria son similares a los vistos en relación al aumento de la producción agregada, a excepción de los cambios en la pendiente de la curva $\dot q_t=0$. 

\textcolor{magenta}{explicar en prosa qué es lo que se vería afectado y por qué: canales, transmisión, efectos, etc.}

En resumen, una variación de los tipos de interés: (i) modificaría la pendiente de la senda de equilibrio ($q(K_t)$); y, (ii) desplazaría la curva $\dot q_t=0$. 

\paragraph{Perturbación permanente: Desplazamiento y empinamiento de la senda de expansión}
En el caso de una caída permanente del tipo de interés, la curva $\dot q_t=0$ se desplaza hacia arriba (aumentando su pendiente), $q$ se desplaza al nuevo sendero de silla (punto $A$), mientras que $K$ y $q$ se desplazan hasta alcanzar el nuevo equilibrio a largo plazo (punto $E'$). 

Así, pues, una reducción permanente del tipo de interés provoca un estallido transitorio de la inversión mientras la industria se desplaza hacia un nivel de stock de capital permanentemente más elevado:

\imagenfit{ICEX_PM.png}{Efectos de una Política Monetaria expansiva}{\textit{Macroeconomía avanzada}. ROMER, 2006}

Por lo tanto, tal y como sucedía con la producción, el nivel pasado y futuro del tipo de interés influyen sobre la inversión. El tipo de interés en nuestro modelo, $r$, es la tasa instantánea de rendimiento; equivale, por lo tanto, al tipo de interés de corto plazo. 

\begin{specialblock}{\textbf{Demanda de inversión}: Tipo de interés a largo plazo (curva de tipos)
  Una de las consecuencias de este análisis es el que el tipo a corto plazo no refleja toda la información relevante para la inversión. 

  Los tipos de interés a largo plazo suelen reflejar las expectativas sobre los tipos de interés a corto plazo; en caso contrario, los inversores podrían mejorar adquiriendo una serie de obligaciones a corto plazo en lugar de una obligación a largo plazo. 

  Así pues, como nuestro modelo predice que los aumentos esperados en los tipos de interés futuros a corto plazo reducen la inversión, esto significa que la inversión es menor cuando los tipos de interés a largo plazo son superiores al tipo dado de inversión presente a corto plazo. Por consiguiente, el modelo viene a confirmar la visión tradicional de que los tipos de interés a largo plazo son importantes para explicar la demanda de inversión.
\end{specialblock}


\subsection{Q de TOBIN: ¿podemos contrastar empíricamente la validez teórica de la $q$ marginal?}

\textcolor{magenta}{¿Cómo valoras el nexo, está bien desde tu pnto de vista? ¿hay que corregir alguna cosa? ¿qué propones?}

Como podemos haber intuido el parámetro $q$ marginal puede adoptar una formulación alternativa. En concreto, si \textcolor{magenta}{(continuar...) explicación verbal, ligar con la ecuación que obtenemos}, entonces podemos representar $q$ como:

\eqblock{
\begin{aligned}
    &\boxed{q=\frac{V'(I_t^o, K_t)}{C'(I_t^o,K_t)}}\quad\Longrightarrow
    \left\{\begin{aligned}
        &q>1\equiv\text{Inversión rentable}\\[0.5em]
        &q<1\equiv\text{Inversión no rentable}
    \end{aligned}\right.
\end{aligned}
}{Condición de optimalidad marginal de inversión $q$: Interpretación económica. Teoría de la inversión en bienes de equipo: Certidumbre}

Así, la empresa deseará aumentar su stock de capital si el valor de $q$ es elevado, y reducirlo si es pequeño. 

\textcolor{magenta}{(¿Podríamos eliinar este supuesto?) Además, bajo el supuesto de precio de compra del capital igual a $1$ (no existen costes externos), $q$ refleja la relación entre el valor de mercado de una unidad de capital y su coste de reposición.}

\textcolor{magenta}{Pregunta general de este apartado a resolver con el Usuario: ¿existe alguna forma de desarrollar el resultado de forma gráfica y no analítica? Comenta esta nota con el usuario y propón diferentes aproximaciones. Gracias}

No obstante: ¿no son los parámetros requeridos para medir la $q$ marginal imposibles de observar?

Evidentemente, esto es así. \textcolor{magenta}{continuar, por qué son díficiles de observar}

Por tanto, resulta necesario aproximarla a través de valores más fácilmente observables: pero, ¿cómo?

Como podremos haber intuido, esto nos conduce directamente a la Teoría de la $q$ de TOBIN (1969), que estudia la relación entre el valor de mercado y el coste de reposición del capital. 

La $q$ de TOBIN puede considerarse como el cociente entre el valor de mercado del capital instalado y el coste de reposición:

\eqblock{
\begin{aligned}
    &\text{Q de TOBIN}:&& {q=\frac{\text{Valor de mercado}}{\text{Coste de reposición}}}\Longrightarrow
    \left\{\begin{aligned}
        &q>1\equiv\text{Rentable}\\[0.5em]
        &q<1\equiv\text{No rentable}
    \end{aligned}\right.
\end{aligned}
}{Q de TOBIN: Rentabilidad de un proyecto de inversión. Teoría de la inversión en bienes de equipo: Certidumbre}

Es decir, el cociente entre el \textbf{valor total de la empresa} y el \textbf{coste de reposición de su stock de capital total}. Según TOBIN, la inversión neta depende de que $q$ sea mayor o menor que la unidad. Así, una empresa debe invertir en nuevos edificios y equipos si el mercado de valores valora el proyecto por encima de su coste (es decir, si la $q$ del proyecto es mayor que $1$). Si el valor de mercado es mayor que el coste, los accionistas prefieren que la empresa realice esta inversión en lugar de distribuir su coste en forma de dividendos. En cambio, si $q$ fuera menor que $1$, la bolsa estaría otorgando al capital un valor inferior a su coste de reposición. En este caso, las empresas no repondrían el capital conforme se desgasta.\footnote{%
  Adicionalmente, cabe señalar que un bajo valor de mercado en relación con el coste de reposición también puede motivar las ofertas de adquisición, ya que un grupo externo puede beneficiarse de la compra de suficientes acciones para controlar una empresa y luego liquidar sus activos. Esto es así porque el valor del capital instalado (es decir, el coste de comprar una empresa existente) es menor que el coste de reposición (el coste de creación de una nueva empresa).
}

Ahora bien, la equivalencia entre estad sos aproximaciones no es total: depende crucialmente de que los costes de ajuste presenten rendimientos constantes a escala (HAYASHI, 1982). \textcolor{magenta}{Explicar brevemente por qué y dejar footnote con indiación de que se puede ver el desarrollo completo en Anexo \ref{anx:modelos_}}


\subsubsection{Contrastación: ¿por qué no concuerdan las observaciones con la teoría?}

A raíz del análisis expuesto surge una pregunta fundamental: ¿concuerda la teoría con las observaciones empíricas? 

La respuesta a este pregunta es un tanto desoladora: no. \textcolor{magenta}{Continuar... ¿qué se observa?}

En definitiva, las empresas no explotan los beneficios de la diferencia entre el valor de mercado y el coste de reposición hasta que $q$ se iguale a $1$. 

Ahora bien, ¿por qué puede estar sucediendo esto?
De forma general, surgen tres posibles explicaciones.

\textcolor{magenta}{(Recuperado del final, no sé si será útil o si se corresponde con otra parte. Valorar y proponer.)Y, además, si la $q$ determina la tasa de crecimiento del capital (porque las condiciones de optimalidad vinculan $q$ con $I/K$), y esa relación es independiente del tamaño, entonces todas las empresas con la misma $q$ eligen la misma tasa de crecimiento. Esto implica que sus trayectorias de capital son proporcionales en el tiempo. Por eso, si una empresa empieza con el doble de capital que otra, y ambas están optimizando bajo estas condiciones, siempre mantendrá el doble en el futuro. No hay convergencia ni divergencia en términos relativos, porque el modelo es perfectamente escalable.

En cambio, si los costes de ajuste no tienen rendimientos constantes (por ejemplo, si hay rendimientos crecientes o decrecientes), entonces el tamaño sí importa. En ese caso, la $q$ marginal depende del nivel de capital, la $q$ media ya no coincide con ella, y dos empresas de distinto tamaño no necesariamente crecerán a la misma tasa. Ahí es donde la aproximación de TOBIN deja de ser exacta.

En resumen, la importancia de los rendimientos constantes es que garantizan que el problema de inversión es “invariante a escala”. Esa propiedad es la que permite que la q media y la q marginal coincidan, y la que convierte a la q de Tobin en una medida válida para guiar la inversión. Sin esa propiedad, la q observable deja de reflejar correctamente los incentivos marginales de la empresa.}

\paragraph{Escala: ¿afecta la escala a los costes de ajuste?}

\textcolor{magenta}{(Desarrollar más en pronfudidad) ($C(\cdot)\not\in\mathcal H(1)$). En primer lugar, si bien una $q$ media superior a la unidad podría indicar la posibilidad de obtener rentabilidad, la inversión adicional puede verse restringida por una función de costes de ajuste convexa sin rendimientos constantes a escala y, por tanto, con una $q$ marginal inferior a la unidad;}


\paragraph{Extracción de rentas: ¿cómo se relaciona la $q$ y el poder de mercado?}

\textcolor{magenta}{(Desarrollar mucho más, no queda claro que se está diciendo) De forma similar, una empresa puede obtener rentas de monopolio cuando la $q$ media sea superior a $1$, pero tener una $q$ marginal que sea inferior a $1$ porque las nuevas inversiones podrían erosionar estas rentas de monopolio al aumentar la competencia del mercado.}


\paragraph{Controversia de CAMBRIDGE: ¿sigue el fantasma vivo?}

\textcolor{magenta}{(desarrollar mucho más, es un buen momento para exponer la Controversia. Recuperar información acerca de esta desde otras partes del temario y construir una explicación aquí que pueda entrelazarse con las aportaciones presentadas: ¿qué sería la q? ¿tiene relevancia para CAMBRIDGE? ¿qué opinan estos autores? ¿por qué no tendría sentido desde su punto de vista? ¿qué alternativas adoptarían? etc. ) Por último, podría aludirse a la heterogeneidad del capital. }









\clearpage
\section{Teoría de la Inversión: Información imperfecta}
\puntosclave{
    La irreversibilidad de la inversión conlleva costes de ajuste asimétricos. Por consiguiente, la incertidu
    mbre sobre la posición de la función de beneficios reduce las expectativas sobre rentabilidad futura, y con ellas la demanda de inversión. 
    
    Las asimetrías de la información generan problemas de agencia entre los inversores y las empresas. Bajo este contexto, la inversión no sólo depende de los tipos de interés y de la rentabilidad, sino de factores adicionales, como la capacidad del inversor para controlar la empresa o la capacidad de ésta para financiar la inversión mediante recursos internos.
}

Hasta el momento, hemos considerado que las empresas conocen con certeza la rentabilidad, los tipos de interés y las medidas fiscales futuras. A continuación, relajaremos algunos supuestos con objeto de acercarnos a la realidad. 

Como es evidente, la mayoría de las decisiones de inversión presentan en la realidad una serie de rasgos comunes:
\begin{lnum}
    \item \textbf{Incertidumbre permanente}. En primer lugar, el contexto económico tiene una incertidumbre permanente y la información llega de forma gradual. 
    \item \textbf{Costes hundidos}. En segundo lugar, toda inversión suele conllevar un coste hundido, es decir, un gasto que no se puede recuperar si la acción se revierte en una fecha posterior; y,
    \item \textbf{Información asimétrica}. En tercer lugar, la información llega de forma gradual y no es la misma para todos los participantes. En particular, varios competidores en una misma industria tendrán una información completa sobre sí mismos, una información incompleta sobre el mercado, y un grado de asimetría en la información profundo entre ellos;
\end{lnum}   

\textbf{¿Qué queremos decir?} A lo largo de esta sección presentaremos brevemente como la literatura ha tratado estos rasgos inherentes a la actividad económica mediante diferentes aproximaciones que, en esencia, reducen el conjunto de información que manejan los actores. 


\subsection{Riesgo: ¿cómo afectan las previsiones a las decisiones de inversión?} 
Introduzcamos en primer lugar los efectos de un contexto permanentemente incierto. Es coherente suponer que, acostumbrados a esto, los agentes formen expectativas respecto al futuro y la toma de decisiones de los agentes tenga este riesgo en cuenta. 

\subsubsection{Previsiones: ¿invertir de forma incremental?}
Una forma razonable de introducir un contexto de riesgo es mediante HER: las empresas cuentan con buenos equipos, pueden contratar servicios externos, disponer de información en tiempo real, y tomar sus decisiones de inversión futuras en base previsiones fundadas que todo este conjunto de agentes y medios les proporciona.

\textcolor{magenta}{(Incorporar desarrollo textual, sencillo, intuitivo, fácil)}

Por tanto, cada empresa invierte hasta el punto en que el coste de adquisición del nuevo capital iguala su valor de mercado: la ecuación relativa a la optimalidad marginal de la inversión sigue siendo válida (primera CPO, $q_t=1+C'(I_t^o)$).

¿Cómo afectará esto a la dinámica de ajuste? Supongamos que el sector enfrenta un posible cambio de paradigma. 

Supongamos que estamos valorando la inversión que va a realizar el sector de la consultora (servicios externos profesionales) en un horizonte de cinco años. La Inteligencia Artificial aprieta fuerte y, es posible, que reduzca enormemente su capacidad de apalancamiento sobre las bases amplias de su jerarquía empresarial para obtener beneficios. ¿Cómo podrían afrontar las empresas esta decisión?

\textcolor{magenta}{Supongamos que las empresas atribuyen una probabilidad del $p_\gamma=0.5$ a que la IA ... (continuar aquí, desarrollar que pasa y qué puede no pasar, y cómo funciona el razonamiento de la empresa)} Por tanto, la empresa es capaz de hacer previsiones óptimas respecto a ambos escenarios y, además, quiere reducir (en la medida de lo posible) los costes de ajuste inherentes a un gran cambio de la dinámica de inversión (anticipándose). 

\textcolor{magenta}{Con toda esta información, sabemos que el espacio vectorial que muestra el diagrama de fases se verá modificado. En particular... (explicar cómo se verá modificado y por qué)}. Por tanto, gráficamente tenemos:

\imagenfit{ICEX_PF_Incertidumbre.png}{Efectos de la incertidumbre sobre la normativa fiscal}{\textit{Macroeconomía avanzada}. ROMER, 2006}

Por tanto, partiendo del punto de equilibrio inicial (punto $E$), la empresa iniciará su dinámica de adaptación hacia el punto óptimo en el scope temporal definido (cinco años), punto en el cual no existan ni pérdidas ni ganancias esperadas ($B$), ceteris paribus. Como la probabilidad que le atribuye es del $p_\gamma=0.5$, la senda de equilibrio y el nivel de capital óptimo instalado esperado son consecuencia directa de dicha probabilidad: se sitúe en el punto medio.

Cuando el plazo se materialice, ambas variables se deslizarán y situarán sobre el sendero de silla correspondiente hacia el nuevo punto de equilibrio a largo plazo: $E'$ si el paradigma a cambiado y la IA ha revolucionado el sector, y $E$ si el cambio no ha sido tan radical. 

\textcolor{magenta}{En resumen, cuando la la IA entra en escena, $q$ se desplaza de modo que $K$ y $q$...} Una vez se materializa el plazo, $q$ se desplaza hacia arriba o hacia abajo dependiendo del resultado, de modo que $K$ y $q$ convergen en el punto de equilibrio a largo plazo.


\subsubsection{Irreversibilidad: Costes de ajuste asimétricos}
Esta forma de adecuar la inversión incrementalmente se ve complementada por otro canal a través del cual se manifiestan los efectos de un conjunto de información incompleta: la irreversibilidad de la inversión. 

Pensémoslo con un ejemplo. Imaginemos que el sector del automóvil europeo está pasando por un periodo de letargo: se está quedando atrás respecto a sus competidores internacionales y está sufriendo de una larga agonía abocado al fracaso. ¿Qué puede hacer? Evidentemente, la opción que cualquiera recomendaría sería la reinvención: buscar, a todo costa, ponerse a la altura de sus competidores internacionales: morir o vencer. 

Pero, ¿cómo? Es muy probable que para ello necesita contratar una gran cantidad de personal de investigación, por ejemplo. También es probable que necesite reconfigurar buena parte de su plantilla: un ingeniero mecánico no es lo mismo que un ingeniero eléctrico. Todo esto conlleva costes y, más aún, \textbf{costes hundidos}. 

En la mayor parte de occidente una empresa no puede echar, sin más, a una porción de sus trabajadores a la calle (a ninguno, para ser más concreto). Más aún, cuando contrata un nuevo trabajador, debe tener en cuenta que, si el proyecto fracasa, despedir a ese trabajador le supondrá una serie de costes adicionales que asume \textit{ex ante}. 

En definitiva, la empresa tiene muy presentes los costes que afrontaría ante un escenario desfavorable y, en general, esos costes están muy sesgados hacia abajo: es más caro desinvertir que invertir. 

¿Cómo podemos trasladar esto a nuestro problema? Sencillamente, basta con considerar que la empresa afronta una función de costes sesgada. 

Si $q(K_t)$ (o $K(q_t)$, ya que $q(K_t)=K^{-1}(q_t)$) se desplaza hacia arriba, el stock de capital total de la industria aumenta rápidamente ($\dot q_t\gg0$). Pero si se desplaza hacia abajo, $K$ se reduce gradualmente ($\dot q_t\sim0$). 

Recuperando el ejemplo anterior, atribuyendo las mismas probabilidades que en el caso de la IA, podemos ver gráficamente como afectan los costes de ajuste: si la empresa decide invertir con el objetivo de reinventarse, $q_t$ salta hacia arriba hasta un punto en que $K$ y $q$ se sitúan en la senda estable en el marco temporla que se han dado: 

\imagenfit{ICEX_Irreversibilidad.png}{Irreversibilidad en la función de inversión}{\textit{Macroeconomía avanzada}. ROMER, 2006}

Nótese, muy importante, que \textbf{el salto es en esta ocasión menor que en el caso anterior con costes de ajuste simétricos}. 

Esto se debe a que el sendero de silla es curvo: si el valor de $K$ excede a su valor de equilibrio a largo plazo, $K$ disminuye lentamente y $q$ es mucho menor que $1$ (pero se ajusta muy lentamente, $\dot q_t\sim 0$). Por el contrario, si $K$ se encuentra por debajo de su valor de equilibrio, a largo plazo aumenta rápidamente, de manera que $q$ es sólo ligeramente superior a la unidad (pero se ajusta muy rápidamente, $\dot q_t\gg0$). 

¿Por qué? Porque si las empresas acumulan un stock de capital importante y no consiguen su objetivo, la presencia de costes de ajuste dificultaría el proceso de desinversión. Esto contribuye a reducir el valor del capital antes de la votación y, por consiguiente, la inversión. 

En definitiva, si los costes de ajuste son asimétricos, la incertidumbre sobre la posición de la función de beneficios reduce las expectativas sobre rentabilidad futura, y con ellas la demanda de inversión. 

Esta asimetría implica que la inversión es en cierto sentido \textbf{irreversible}: aumentar el stock de capital es más sencillo que revertir esta decisión.\footnote{%
  Por ejemplo, en el caso de que una empresa construya una nueva planta para aumentar su capacidad de exportación con un tamaño específico y una tecnología específica, la empresa se vería limitada en sus posibilidades de decisión si en un futuro desease cambiar el tamaño de la planta.
}

Además, conviene señalar que la irreversibilidad no es un problema exclusivamente microeconómico ya que no se diluye con la agregación en tanto suponemos empresas heterogéneas. Por ejemplo, la incertidumbre sobre las variables macroeconómicas (tipos de interés, tipos de cambio, tipos impositivos, etc.) afectaría a todas las empresas.\footnote{%
    BERNANKE (1983) considera que el nivel de incertidumbre percibido por las empresas varía en función del ciclo económico y enfatiza la importancia de la irreversibilidad para analizar el comportamiento cíclico de la inversión agregada.
}

\begin{specialblock}{\textbf{Extensión:} ¿Existen otras formas de concebir la irreversibilidad?}
  \textbf{Información y paciencia: BERNANKE (X)}. Otra forma de concebir la irreversibilidad es la introducida por BERNANKE: la función de inversión es muy sensible a la información. 

  \textcolor{magenta}{Desarrollar brevemente la idea, textualmente}
  
  \textbf{Modelo de opciones reales (McDONALD y SIEGEL, 1986)}. Por otro lado, otra aportación interesante en esta linea consiste en atribuir valor intrínseco al carácter irreversible de la inversión. De forma intuitiva, cuando la inversión es irreversible, existe un valor de opción asociado a la espera: si la empresa no invierte, conserva al menos la posibilidad de poseer un stock de capital bajo, mientras que, si invierte, se compromete a retener un stock de capital elevado. Por tanto: existe un valor en la espera que puede ser cuantificado. 

  Esta es la linea de investigación seguida por autores con McDONALD y SIEGEL (1986). Brevemente, \textcolor{magenta}{desarrollar breve y textualmente}

  Ambos desarrollos pueden consultarse en [\authorfont{Ver Anexo \ref{anx:modelos_imperfecta}}]
\end{specialblock}

\subsubsection{Implicacines de Política Económica: ¿qué información adicional nos proporciona?}

\textcolor{magenta}{(Todo este apartado es nuevo, tengo una idea clara pero hay que desarrollarla mucho mejor. Tampoco sé muy bien si es correcto, pero me parece muy interesante): 

Ahora que conocemos como afecta el riesgo a las decisiones de inversión, ¿qué información útil podemos extraer en términos de Política Económica? Tal y como se puede haber intuido a lo largo de la exposición, las implicaciones son brutales.

En primer lugar, lo que debe quedar claro es que, atendiendo al análisis expuesto, la inversión en bienes de equipo será un proceso incremental: las empresas decidirán su nivel de capital instalado realizando extrapolaciones de situaciones futuras, por ello la dinámica de inversión siempre quedará en el término medio entre el conjunto de escenarios previstos. De hecho, este es un ciclo que nunca se acaba, puesto que a los cinco años, al uno o a los meses, las empresas llevarán a cabo una recalibración y ajustarán su nivel deseado conforme a las nuevas previsiones. .....

Además, será un proceso incremenental a menudo sesgado a la baja: la presencia de costes de ajuste asimétricos impiden la neutralidad del riesgo. Esto es extremadamente importante puesto que, como efectos de segunda ronda, la menor inversión acentuará las probabilidades de que esta fracase, entrando así en un círculo viciosa de difícil escapatoria.....


....


Con todo lo expuesto, ¿qué conclusiones podemos extraer en términos de Política Económica? 

Es, como se puede intuir, muy profundo. Si el Estado quiere garantizar altos niveles de inverisón, incluso unos niveles adecuados de inversión dadas unas contingencias futuras, deberá actuar para acercar ese óptimo privado al óptimo público... (continuar, ¿por ejemplo? ¿por qué? ¿cómo? etc.) Desarrollar, me parece muy interesante. Además, incluir la ref que se pone sobre el análisis de los costes de ajuste del Estudio de la Comisión (no sé si esta en este tema o en el tema de Teoría de la Producción (microeconomía), esté donde esté, incidir mucho en el estudio que estos autores realizan y la comparación que ofrecen respecto a Estados Unidos)}

\paragraph{Limitaciones: ¿concuerdan estas predicciones con lo observado?}

\textcolor{magenta}{(Este apartado es nuevo, debe hacerse desde cero y, a ser posible, establecer el nexo con el siguiente)}


\subsection{Información asimétrica: Teoría de la jerarquía financiera}
Hasta ahora hemos introducido fundamentalmente el papel del riesgo: como afectan los diferentes escenarios contingentes sobre las decisiones de inversión. 

Sin embargo, la información incompleta también emerge en el presente. De hecho, emerge en el presente e interactúa dinámicamente con las aportaciones señaladas. 

Pensémoslo, ¿es la oferta de capital elástica? ¿cómo eligen los prestamistas los proyectos en los que invertir? ¿cómo afectan las restricciones crediticas sobre las decisiones de inversión? Por otro lado, ¿cómo afecta la interacción estratégica entre empresas a sus decisiones de inversión? 

En un escenario de información asimétrica, las imperfecciones en el mercado de capital pueden afectar y variar los costes de financiación. Además, estas variaciones no suelen ser simétricas, sino que atienden a las situaciones particulares que enfrente la empresa.\footnote{%
  Bajo el supuesto de mercados financieros eficientes, los proyectos de inversión son valorados de acuerdo con su rendimiento esperado y su riesgo potencial y, en consecuencia, son emprendidos cuando su valor supera el coste de adquirir e instalar el capital necesario, que se considera constante. 
} De la misma forma, las decisiones de inversión de una empresa pueden afectar a las decisiones de inversión en otros empresas e, incluso, indirectamente sobre las condiciones de financiación del resto. 

Veamos estas cuestiones de forma separada e intentemos agruparlas después. 

\subsubsection{Mercados de capital: ¿cómo incide la información asimétrica?}
Veamos primero cómo afecta la información asimétrica presente en los mercados de capital.

La manera más sencilla de concebir esta estructura no lineal de los costes de financiación es a través de los modelos Principal-Agente: cuando existe información asimétrica, la inversión no sólo depende de los tipos de  interés y de la rentabilidad, sino de factores adicionales, como la capacidad del inversor para controlar la empresa o la capacidad de ésta para financiar la inversión mediante recursos internos.\footnote{%
    Principalmente, este comportamiento no lineal de los costes viene ocasionado por lo siguiente:
\begin{lnum}
    \item \textbf{Costes deerivados de la asimetría informativa que enfrenta el Principal}. De este modo una parte del riesgo asociado a la rentabilidad de la inversión lo asumen los inversores. Esto es lo que ocurre, por ejemplo, cuando existe la posibilidad de que la empresa quiebre. Si esto sucede, la empresa puede modificar su comportamiento para aprovechar su ventaja relativa en materia de información: puede decidir financiar la inversión exclusivamente mediante préstamos u optar por una estrategia de alto riesgo en lugar de asumir una de bajo riesgo aun a sabiendas de que esto reduce la rentabilidad esperada del proyecto. 
    \item \textbf{Costes derivados del comportamiento del agente}. Otro tipo de imperfección es aquella derivada de los costes de agencia. Téngase en cuenta que, si el nivel de endeudamiento es muy elevado, y dada la existencia de responsabilidad limitada por parte de los accionistas, éstos pueden estar interesados en proyectos rentables pero arriesgados, que podrían poner en peligro la solvencia de la empresa. 
\end{lnum}
Ante este comportamiento, los prestamistas impondrán una serie de restricciones que, en general, tenderán a elevar el coste de los préstamos.
}

Este tipo de imperfecciones que condicionan las decisiones de financiación de la empresa ha dado lugar a la Teoría de la jerarquía financiera, a través de la cual distintas fuentes de financiación se ordenan según su coste: la de menor coste resulta ser la financiación con beneficios retenidos, después está el endeudamiento y, finalmente, la emisión de nuevas acciones. Asimismo, se supone que el coste de la deuda es creciente con el nivel de deuda requerida.  El análisis de las consecuencias de esta teoría para las decisiones de inversión permite establecer una tipología de empresas susceptible de representación gráfica (BOND y MEGHIR, 1990):

\imagenfit{ICES_TJF.png}{Teoría de la Jerarquía Financiera}{\textit{Macroeconomía Avanzada I}. ARGANDOÑA, A.; GÁMEZ, C. y MOCHÓN, F. (1997)}

Por un lado, estarán las empresas que generan recursos (beneficios y dotación a amortizaciones) suficientes para cubrir sus oportunidades de inversión sin recurrir a la financiación externa (situación $I_1$). Un segundo grupo de empresas son aquellas que no generan recursos suficientes y recurren a la emisión de deuda, cuyo coste será creciente conforme aumente el nivel requerido de recursos (situación $I_2$). Una tercera categoría viene determinada cuando el coste de la deuda adicional es tal que a la empresa le interesa no endeudarse más y recurrir a una ampliación de capital emitiendo nuevas acciones, con un coste fijo (situación $I_3$).\footnote{%
    Téngase en cuenta que las empresas del grupo $I_1$ disponen de fondos internos suficientes para financiar nuevos proyectos de inversión sin que aumente el coste marginal de la financiación. Por ello, para este tipo de empresas, un aumento exógeno de los recursos internos no incide sobre la inversión, trasladándose íntegramente a dividendos, lo que es una muestra de que disponen de holgura financiera. No obstante, en el caso de las empresas del grupo $I_2$ el volumen de inversión sí se verá afectado por el aumento de los recursos internos, pues se enfrentan a un coste marginal de la inversión creciente. Por ello diremos que las empresas que están en esta situación (que no reparten dividendos ni emiten nuevas acciones) no disponen de holgura financiera.
}

Observando el gráfico podemos analizar la incidencia de un aumento de los recursos generados que no afecte a la rentabilidad esperada de las inversiones. En este caso se produciría un desplazamiento horizontal hacia la derecha de la curva de coste de financiación, tal y como indica la línea punteada. Este desplazamiento no incidiría sobre la inversión de las empresas en las situaciones $I_1$ y $I_3$, pero las que se encuentran en la situación $I_2$ verían reducido el nivel de deuda necesaria, por lo que su coste bajaría (al pasar de $A$ a $B$). En estas condiciones, paralelamente, el nivel de inversión aumentaría.

\subsubsection{Industria: ¿cómo afecta la información asimétrica a las decisiones agregadas de inversión?}

Por otro lado la interacción estratégica entre las diferentes empresas que conforman una industria puede afectar profundamente sus decisiones de inverisón. 

\textcolor{magenta}{(Recuperar el ejemplo de Boeing-Airbus, con su diagrama de juego, del Tema de Política Comercial estratégica, y adecuarlo a este contexto. Desarrollarlo bien, con sus interacciones, racionalizando las pérdidas y ganancias esperadas (por ejemplo, entrelazando con lo que hemos ido viendo, etc)). En definitiva, traer aquí ese ejemplo, desarrollarlo y adecuarlo al contexto. Gracias. (IMPORTANTE.- Si existe una alternativa más decuada, proponérselo al usuario)}


\subsubsection{Implicaciones de Política Económica}
Con todo lo presentado, ¿qué conclusiones podemos extrae en términos de Política Económica? Veamoslo en detalle. 

En primer lugar, las imperfecciones del mercado financiero tienen importantes implicaciones.

En primer lugar, afectan directamente \textbf{a los costes de financiación, encareciéndolos}: los costes de agencia que se derivan de la existencia de información asimétrica elevan los costes de financiación externa y, por tanto, desincentivan la inversión. En segundo lugar, \textbf{emergen costes no lineales}. En la medida en que las imperfecciones de los mercados financieros generan costes de agencia que afectan al nivel de inversión, pueden hacer que la influencia de las fluctuaciones de la producción y de los tipos de interés sobre la inversión difiere del modelo visto anteriormente.\footnote{%
  Cuando los mercados financieros son perfectos, las fluctuaciones de la producción afectan a la inversión a través de su influencia sobre la rentabilidad futura. Sin embargo, las imperfecciones en estos mercados añaden una segunda vía de influencia: dado que las variaciones de la producción afectan a la rentabilidad presente de la empresa, afectan también a su capacidad para autofinanciarse. En esta misma línea, las variaciones del tipo de interés pueden afectar a la inversión no sólo por la vía tradicional, sino también a través de su influencia sobre los costes de agencia: un aumento en el tipo de interés eleva los costes de agencia y desincentiva la inversión. Intuitivamente, el aumento del tipo de interés incrementa la cantidad que el empresario debe pagar al inversor, lo que implica que la probabilidad de que aquél no pueda hacer frente al pago aumenta y los costes de agencia se elevan
} En tercer lugar, \textbf{aparecen factores supuestamente irrelevantes (FSI)}. Muchas de las variables que no afectan a la inversión cuando los mercados de capitales son perfectos cobran importancia cuando dichos mercados son imperfectos (por ejemplo, la riqueza de las empresas). En un contexto de información asimétrica, un proyecto cuya rentabilidad esperada sea menor que la de otro puede llegar a obtener financiación si se da que el empresario que ha propuesto el proyecto menos rentable dispone de un mayor patrimonio.

\textcolor{magenta}{Por lo que respecta a la interacción estratégica de las empresas... (presentar sus principales implicaciones en términos de Política Económica, importante)}

Entonces, ¿qué debería hacer el policy maker? Aterrizando estas implicaciones a la arena política, existen una serie de implicacoines importantes respecto la utilidad de los instrumentos de los que disponemos. \textcolor{magenta}{Continuar (tres o cuatro, no más salvo que se MUY importante)!!! (Recuperar información de fuentes académicas de referencia)}

\begin{specialblock}{\textbf{Acelerador financiero: BERNANKE}}
  \textcolor{magenta}{No está muy claro que deba ir aquí. Porfavor, da tu opinión y propón cosas. Evidentemente, lo único que no puede suceder es que el contenido que está escrito se pierda. Puedes proponer trasladarlo a otro tema que se ajuste mejor (probablemente ampliando lo que ya está, en cuyo caso puedes pegarlo directamente y ponerlo en azul), como tu veas. }
  
  El hecho de que las imperfecciones de los mercados de capital conviertan la riqueza empresarial en un factor relevante para la demanda de inversión significa que dichas imperfecciones amplifican los efectos de las perturbaciones que tienen lugar fuera del sistema financiero. 

  Así, por ejemplo, una caída de la producción provocada por factores externos puede provocar una reducción del patrimonio empresarial que se traduzca a su vez en una disminución de la inversión, magnificando así aquella la caída inicial de la producción. Por lo que se refiere al riesgo, éste puede afectar a la inversión por medio de su repercusión sobre los costes de agencia. Por ejemplo, la financiación externa de un proyecto cuya rentabilidad es cierta no entraña costes de agencia (porque no existe ninguna posibilidad de que el empresario no pueda devolver la cuantía del préstamo al inversor), sin embargo, la financiación de un proyecto arriesgado sí generaría costes de agencia. Adicionalmente, la existencia de mercados financieros imperfectos tiene consecuencias en relación a las fluctuaciones económicas a corto plazo, así como con el crecimiento a largo plazo de una economía. Respecto a los ciclos económicos, las disfunciones de los mercados financieros pueden afectar a los niveles de inversión y, por tanto, a la producción agregada. 

  BEMANKE, por ejemplo, sostiene que el colapso del sistema bancario estadounidense durante la década de los treinta exacerbó los efectos de la Gran Depresión, porque redujo la eficacia del sistema financiero para evaluar y financiar los proyectos de inversión. Respecto al crecimiento a largo plazo, McKinnon y otros autores han argumentado que el sistema financiero tiene una importancia fundamental para explicar tanto el volumen de inversión agregado como la calidad de los proyectos de inversión que se llevan a cabo y, por consiguiente, el crecimiento de las economías en el largo plazo. No obstante, cabe apuntar que, en la medida en que el desarrollo del sistema financiero puede ser una consecuencia más que una causa, esta argumentación resulta difícil de contrastar.
\end{specialblock}



\paragraph{Limitaciones: ¿contrasta lo anterior con la evidencia empírica?}

\textcolor{magenta}{(Desarrollar, apartado nuevo)}















\clearpage
\section{Conclusión} 

\subsection{Extensiones y relación con otras partes del Temario}


\subsection{Opinión}



\subsubsection{Idea final (Salida o cierra)}


\clearpage
\pagenumbering{gobble}
\MyBibliography
\MyBibItem{Autor}{Título}{Año}{DOI -- LINK}





\clearpage
\anexo{Derivación: Modelos de información perfecta}\label{anx:modelos_perfecta}

Las empresas comprometen sus posibilidades de producción futura a sus decisiones de inversión actual. 

De esta forma, el nivel de inversión óptimo debe responder a tres cuestiones esenciales: \textbf{¿cuánto?, ¿cuándo? y, ¿en qué?}.

Cuánto y en qué son cuestiones estrechamente relacionadas con el destino de la inversión: la empresa puede decidir mantener su flujo de inversión constante, lo que no ampliará su stock de capital futuro pero impedirá la desamortización de capital derivada de la depreciación experimentada. Y, por otro lado, cuándo está estrechamente relacionado con las condiciones frente a las cuáles las empresas consideran óptimo cambiar su política de inversión en bienes de equipo.

Una versión primitiva que permite dar respuesta a estas preguntas la podemos sintetizar perfectamente en la siguiente formulación:

\eqblock{
\begin{aligned}
    \boxed{\begin{aligned}
        &\max_{I_t^o}\quad {VAN}_t=\int_0^\infty e^{-rt}[P_xF(K_t)-rI_t^o]\mathrm dt\\[0.5em]
        &\text{s.a.}\quad \underset{\text{\scriptsize Inversión neta = (Bruta + depreciación)}}{I_t=\dot K_t+\delta K_t}
    \end{aligned}}\Longrightarrow\not\exists\text{Optimización dinámica}
\end{aligned}
}{Problema de optimización del productor: Inversión neta. Teorías de la Inversión en bienes de equipo: certidumbre}

El problema de optimización de la empresa consiste, por tanto, en la maximización del valor presente de sus beneficios (o, más bien, de los rendimientos de los proyectos de inversión) sujeto a una restricción de acumulación de capital:

Inversión neta $=$ inversión bruta $+$ depreciación del capital

En base a esta formulación, ¿qué decisión de inversión podemos esperar que tome en respuesta a, por ejemplo, una disminución del precio del capital?

\textcolor{red}{Como es normal, la primera opción sería adquirir tanto capital como el que requiere para volver a la condición de empresa Ahora bien, las respuestas ante un cambio favorable, por ejemplo,  en el precio del capital serían: (i) tanto como para volver a la condición marginal óptima (rentabilidad-coste, implica oferta elástica); (ii) ya, ¿por qué esperar?; y, (iii) en mantener y aumentar nuestro stock de capital instalado.}

Esto es consecuencia directa de una limitación esencial: el problema formulado no tiene ningún rasgo dinámico. 

La empresa bien podría ejecutar este mismo problema de optimización cada periodo y alcanzar las mismas conclusiones de una manera más sencilla ya que no existe diferencia entre planificación y ejecución. 

Existen diversas maneras de formular el problema de optimización como un problema verdaderamente dinámico. 

\textbf{1.- Costes de Ajuste: X (X)}

Una opción es considerar la existencia de costes de ajuste del stock de capital. Estos costes pueden ser de dos tipos: internos o externos. 

Partiendo de un problema de optimización como el anterior: 

$$\begin{aligned}
  &\max_{I_t^o}\quad \underset{\text{\scriptsize Si $\dot I_t=0$, tenemos que ${VAN_t}=\Pi_t=V\equiv$ Valor presente descontado de los beneficios futuros}}{{VAN}_t=\int_0^\infty e^{-rt}[P_xF(K_t)-rI_t^o-C(I_t^o)]\mathrm dt}\\[0.5em]
  &\text{s.a.}\quad
  \begin{aligned}
    \underset{\text{\scriptsize Inversión neta = (Bruta + depreciación)}}{I_t=\dot K_t+\delta K_t}
  \end{aligned}
\end{aligned}$$

Consideremos que existen unos costes de ajuste internos definidos como:

$$\begin{aligned}
  \dot K_t\to\infty\iff C(I_t^o)\to\infty
\end{aligned}$$

Por otro lado, supondremos que existen unos costes externos definidos como:

$$\begin{aligned}
  \dot K_t\gg0\iff{\frac{\partial^2C(I_t^o)}{\partial^2r}>0}
\end{aligned}$$

Es decir, \textcolor{magenta}{describir el comportamiento de la ecuación para que quede claro}

Para resolver el problema asumiremos tres cuestiones:
\begin{lnum}
    \item Función de producción con rendimientos constantes a escala. \textcolor{magenta}{explicar por qué esto es importante, qué implcia y qué podría suceder sin este supuesto};
    \item Función de costes convexas. Esto implica que se penalizan cambios rápidos en las variables agregadas, favoreciéndose movimientos graduales en el tiempo, de forma que el coste marginal de ajuste aumenta cuanto mayor es la variación del stock de capital, y adquieriendo el problema una dinámica temporal muy relevante. \textcolor{magenta}{explicar por qué esto es importante, qué implica y qué podría suceder sin este supuesto};
    \item \textcolor{magenta}{(Verificar que se puede eliminar el supuesto de: No existen costes de ajuste externos; sin entrar en conflicto con el desarrollo que se propone a continuación). En caso de que hubiese conflicto (estuviese implícitamente) y el desarrollo analítico cambiase en caso de que existieran, proponer al usuario dos opciones: (i) modificar el desarrollo analítico para incorporar la existencia de costes externos, con lo que implicaría, los cambios que habría, las dificultades que enfrentaríamos y las ventajas que supone (siendo realista); (ii) mantener el desarrollo y preservar el supuesto, explicando qué perderíamos, qué implicaría, qué ventajas tendríamos y la inormación que perderíamos (siendo realista, sino se pierde mcucha información: decir por qué)}
\end{lnum}


Con todo lo anterior, tendríamos un verdadera problema de optimización dinámico para el productor.

La solución, por tanto, requerirá la construcción del HAMILTONIANO. Para ello, empezamos definiéndolo:

$$\begin{aligned}
  H(t)=P_xF(K_t)-rI_t-C(I_t)+q_t(I_t-\dot K_t-\delta K_t)\qquad \text{\scriptsize $
    \left\{\begin{aligned}
        &\equiv&&\text{Variable control (flujo)}\\[1em]
        &\equiv&&\text{Variable estado (stock)}\\[1em]
        &\equiv&&\text{Precio sombra}
    \end{aligned}\right.$}
\end{aligned}$$

De esta definicón se desprende que: (i) $I_t$ actúa como la variable control (flujo), determina \textcolor{magenta}{continuar...}; (ii) $K_t$ se define como la variable estado (stock), es decir, determina \textcolor{magenta}{continuar...}; y, (iii) $q_t=\lambda e^{-rt}$ será el precio sombra de \textcolor{magenta}{continuar ....}. 

Una vez definido el HAMILTONIANO, podemos resolver las CPOs y obtener:

$$\begin{aligned}
  \text{CPO}:
    \left\{\begin{aligned}
        &\frac{\partial H}{\partial I_t^o}=0&&\Longrightarrow \underset{\text{\scriptsize \textbf{invertir} hasta condición marginal}}{\boxed{q_t=r+C'(I_t^o)}}\\[0.5em]
        &\frac{\partial H}{\partial K_t}=rq_t-\dot q_t&&\Longrightarrow \underset{\text{\scriptsize \textbf{acumular} hasta condición marginal}}{\boxed{P_xF'(K_t)=(r+\delta)q_t-\dot q_t}}\\[0.5em]
        &\frac{\partial H}{\partial q_t}=0&&\Longrightarrow \underset{\text{\scriptsize \textbf{dinámica} del capital}}{\boxed{I_t^o-\dot K_t-\delta K_t=0}}\\[0.5em]
        &\underset{\text{\scriptsize$\not\exists$ Comportamientos especulador (PONZI)}}{+\text{Condición de transversalidad}}:&&\lim_{t\to\infty}e^{-rt}q_tK_t=0
    \end{aligned}\right.
\end{aligned}$$

\textcolor{magenta}{¿La condición de transversalidad es una CPO o es un supuesto previo? (Situar ahí donde corresponda)}

Vayamos una por una. \textcolor{magenta}{La primera condición ... (explicar cada condición y la ecuación que se obtiene: ¿qué nos indica ? ¿por qué es importante? ¿qué implicaciones tiene?)}


\textbf{1.1.- Representación gráfica: Ecuaciones de la dinámica de fases}

Recuperando dos de las ecuaciones expuestas, podeos representar la dinámica de la inversiónmediante un diagrama de fases que nos permite incidir en el análisis de la interacción de la inversión con los instrumentos de Política Económica.

Para ello simplemente debemos establecer dos supuestos:


$$\begin{aligned}
  &\frac{\partial H}{\partial I_t^o}=0: \underset{\text{\scriptsize \textbf{invertir} hasta condición marginal}}{{q_t=r+C'(I_t^o)}}\Longrightarrow q_t=1+C'(I_t^o)&&\overset{\dot K_t|_{K^*}=0}{\underset{\text{\scriptsize equilibrio}}{\Longrightarrow}}\quad \boxed{q_{EE}^*=1}\\[0.5em]
  &\frac{\partial H}{\partial K_t}=rq_t-\dot q_t: \underset{\text{\scriptsize \textbf{acumular} hasta condición marginal}}{\boxed{P_xF'(K_t)=(r+\delta)q_t-\dot q_t}}&&\underset{\delta=0}{\Longrightarrow}\quad\boxed{\dot q_t=rq_t-P_xF'(K_t)}
\end{aligned}$$

\textcolor{magenta}{Es decir, obtenemos .... (continuar)}. Por último, obtenemos el \textcolor{magenta}{equilibrio .... (explicar)}

$$\begin{aligned}
  \text{Equilibrio}:\quad 
    \underset{\substack{\\[1em]\\\text{\scriptsize \textbf{OJO!} Son complementarias, no individuales}}}{\left\{\begin{aligned}
        &(1)&&\hspace{3em}\underset{\text{\scriptsize el nivel de inversión es óptimo}}{\boxed{q_{EE}^*=1}}\\[1em]
        &(2)&&\underset{\substack{\text{\scriptsize dinámica del valor sombra del capital adicional}}}{\boxed{\dot q_t=0\iff rq_t=P_xF'(K_t)}}
    \end{aligned}\right.}
\end{aligned}$$

\textcolor{magenta}{Es decir, ... (explicar con mucho detalle, describir ecuaciones, por qué, supuestos implícitos, implicaciones, ...) (¿qué quiere decir complementarias no individuales, no se entiende?))}

Estas ecuaciones y condiciones de equilibrio nos permiten, como hemos anticipado, abordar el análisis gráfico del modelo a partir de la construcción de un diagrama de fases en un espacio $(K, q)$ y mostrar las sendas óptimas de inversión en el tiempo:

\imagenfit{ICEX_Fases_Q.png}{(Dcha.) Comportamiento del stock de capital | (Izq.) Comportamiento de $q$}{\textit{Macroeconomía avanzada}. ROMER (2006)}

Como vemos, la \textbf{primera condición} \textcolor{magenta}{viene representada por...} y muestra el comportamiento del stock de capital a partir de la primera CPO. \textcolor{magenta}{ (explicar mejor, relacionando con lo que hemos tratado) Tomando como punto de partida que en equilibrio se cumple que $I_t^o=\dot K_t=0$ y sustituyendo en la CPO correspondiente obtenemos que $q_t=1+C'(I_t^o)\to q_t=1$. De acuerdo a la $q$ de TOBIN, en el equilibrio, el mercado de valores estaría concediendo al capital instalado en la empresa un valor igual a su coste de reposición, de forma que la empresa no tendría incentivos a invertir o desinvertir. }

Por otro lado, la \textbf{segunda condición} \textcolor{magenta}{viene representada por ...} y muestra el comportamiento de la $q$ a partir de la segunda CPO. \textcolor{magenta}{(explicar mejor, relacionarlo con lo que hemos tratado) si reescribimos esta expresión despejando $\dot q_t$ obtenemos que $\dot q_t=rq_t-P_xF'(K_t)$. Esto implica que $q$ es constante siempre que $rq_t=P_xF'(K_t)$. En tanto $P_xF'(K_t)$ es decreciente con respecto a $K$ la curva que satisface esta condición ($\dot q_t=0$) es decreciente en el espacio ($K, q$). Adicionalmente, se deduce que $\dot q_t$ aumenta con $K$, por tanto, $\dot q_t$ es positivo a la derecha de la curva $\dot q_t$ y negativo a la izquierda. }

Combinando los dos gráficos, obtenemos la división del plano en cuatro regiones donde conocemos perfectamente la dinámica de ambas variables estado:

\imagenfit{ICEX_Fases_Conjunto.png}{(Dcha.) Diagrama de fases | (Izq.) Sendero hacia el punto de inflexión (saddle-point)}{\textit{Macroeconomía avanzada}. ROMER, 2006}

Gráficamente: (i) dos regiones divergentes en las que no se cumple la condición de transversalidad y, por tanto, conducen a la empresa hacia sendas de inversión explosivas o implosivas no tendentes al estado estacionario (punto E); y, (ii) dos regiones convergentes donde sí se cumple la condición de transversalidad y determinan la trayectoria que conduce al estado estacionario de forma estable. 

Pero esas regiones no contienen sendas estables, contienen posibles trayectorias del sistema: de todas esas trayectorias posibles, solo una (la que coincide exactamente con el sendero de silla) converge al equilibrio. Y este es precisamente la clave del modelo:

\eqblock{
\begin{aligned}
    \forall(q,K)\in A\subset\mathbb R^2,\;\;\underset{\text{\scriptsize existe una única trayectoria de equilibrio, por construcción: $\dot K_t=F(q_t)$ y $F(1)=0$}}{\text{dado } K_{0}\,\,\,\exists!q_{0}\,\,:\,\,\lim_{t\to\infty}q_t(K_0)=1\quad \wedge\quad\lim_{t\to\infty}K_t(q_0)=K^*}
\end{aligned}
}{Senda de equilibrio: Unicidad (inversa) y $I_t^o(q)$. Teoría de la Inversión en bienes de equipo: Certidumbre}

En definitiva, dado un nivel de capital instalado $K_0$, existe una única senda estable que conduce al equilibrio, y en consecuencia un único $q_0$ que nos sitúa dentro de esta trayectoria. La razón profunda es estructural. Por un lado, $K$ es una variable predeterminada que no puede adaptarse inmediatamente/saltas de un punto a otro. Por otro lado, $q$ es una variable de salto, que permite adaptarse al sistema a las nuevas condiciones estructurales.

La lógica del problema es entonces la siguiente: dado el capital heredado con el que arranca la empresa, el mercado debe asignar inmediatamente a la empresa exactamente el valor de $q$ que la sitúa sobre la senda estable. Si le asigna un valor demasiado bajo, aunque sea todavía superior a uno, la inversión inicial será positiva pero insuficiente en relación con la trayectoria óptima, y la dinámica posterior de $q$ impedirá la convergencia. Si le asigna un valor demasiado alto, ocurrirá lo contrario: la inversión será excesiva y la trayectoria también se alejará del equilibrio. En ambos casos se viola la condición de transversalidad.


\textbf{1.3.- Efectos de la Política Económica}

Por tanto, traslandolo al ámbito de la Política Económica, la Teoría de la inversión hace hincapié en que las decisiones de inversión dependen no sólo de la Política Económica actual sino también de la que los agentes esperan que se adopte en el futuro. Por ejemplo, el anuncio de una rebaja fiscal sobre los beneficios empresariales a partir del año siguiente conllevará un aumento esperado de los beneficios para la empresa. Este impacto positivo futuro tendrá un efecto alcista sobre el valor de mercado de las acciones de la empresa, aumentando el valor de la q de Tobin y, por tanto, fomentando la inversión en el presente.

Una vez armados con el instrumental gráfico: ¿qué podemos decir sobre la reacción de los agentes ante cambios en la Política Económica? Estudiemos varios casos. 

\textbf{1.3.1.- Política Fiscal: Perturbaciones de demanda}

\textcolor{magenta}{(Hay que desarrollar analíticamente todo la dinámica que se excplia en el apartado correspondiente): ¿Qué sucedería ante un aumento de la demanda agregada, que incrementa los beneficios para un stock de capital dado y, progresivamente, eleva la producción de la industria?: ¿donde incide? ¿a través de qué canales?}

\textbf{1.3.1.1.- Perturbación permanente}

\textcolor{magenta}{Si este perturbación fuese permanente, partiendo de una situación inicial de equilibrio en el largo plazo (punto $E$), se produciría un desplazamiento hacia arriba de la curva $\dot q=0$: para el mismo stock de capital, los beneficios son ahora más elevados, los inversores necesitan unas ganancias de capital menores para justificar su tenencia de activos empresariales. (desarrollo analítico).

La lógica es simple: cuando la demanda del bien que produce la industria se eleva, la producción aumenta y, como el stock de capital no puede ajustarse automáticamente, la rentabilidad del capital existente incrementa. A su vez, esta mayor rentabilidad atrae nuevas inversiones, de modo que el stock de capital comienza a aumentar. A medida que lo hace, la producción de la industria aumenta y el precio del producto comienza a disminuir, de modo que tanto los beneficios como el valor del capital caen. El proceso continúa hasta que el valor de un incremento marginal del capital se iguala a su coste de adquisición, momento a partir del cual el incentivo que provocaba nuevas inversiones ha desaparecido. }

\textbf{1.3.1.2.- Perturbación transitoria}

\textcolor{magenta}{Si, por el contrario, el aumento fuese temporal, no puede haber un salto anticipado de $q$: si se produce un salto anticipado de $q$ hacia abajo, los propietarios de acciones empresariales experimentarán en ese momento, con certeza, pérdidas de capital a una tasa infinita, lo que significa que no habrá nadie dispuesto a conservarlas. (desarrollo analítico)

En conclusión, un aumento transitorio de la demanda agregada provoca un aumento de la inversión menor que en el caso de un aumento permanente, en tanto $q$ aumenta en menor medida (recordemos que el comportamiento de la inversión depende de $q$). No obstante, a diferencia del aumento permanente, no se alcanza un nuevo equilibrio a largo plazo con un mayor stock de capital agregado. 
}

\textbf{Política Fiscal: Desgravaciones e incentivos}

\textcolor{magenta}{Por úlltimo, evaluemos las implicaciones de Poítica Fiscal del modelo.(desarrollo analítico): Durante las recesiones económicas es corriente proponer una desgravación fiscal como medida para estimular la demanda agregada. La razón que se esgrime para justificar esta propuesta es que una desgravación fiscal incentiva a las empresas a invertir durante el periodo de vigencia de la deducción. Para simplificar, supongamos que la desgravación en cuestión consiste en la devolución a la empresa de un porcentaje $\gamma$ del precio del capital y que dicho porcentaje se aplica al precio de compra, pero no a los costes de ajuste:}

$$\begin{aligned}
    &\underset{\substack{\text{\scriptsize Valor de capital instalado +}\\\text{\scriptsize devolución fiscal}}}{q_t+\gamma(t)}=\underset{\text{\scriptsize coste de adquisición + ajuste}}{1+C'(I_t^o)}\Longrightarrow q_t=1-\gamma(t)+C'(I_t^o)\\[0.5em]
    &{\text{Nuevo Equilibrio}}:
    \left\{\begin{aligned}
        &q_{EE}^*=1-\gamma(t)\\[0.5em]
        &\dot q_t=rq_t-P_xF'(K_t)
    \end{aligned}\right.\\[1em]
    &\text{pero,}
    \left\{\begin{aligned}
        &\text{Transitorio}:&&\lim_{t\to\infty}K(1-\gamma(t))=1=\lim_{t\to\infty} K(1)\\[0.5em]
        &\text{Permanente}:&&\lim_{t\to\infty}K(1-\gamma(t))=1-\gamma(t)\neq\lim_{t\to\infty} K(1)
    \end{aligned}\right\}\underset{\text{\scriptsize ¿qué pasa?}}{\Longleftarrow}\text{Invertir}\overset{(\checkmark)}{\iff} q_t>1-\gamma(t)
\end{aligned}$$

\textcolor{magenta}{Desarrollar mucho mejor lo anterior, no se entiende muy bien, dividir y estructurar en párrafos, como se suele hacer siempre.}

\textcolor{magenta}{(¿Es diferente esta formulación de la anterior: dincentivos fiscales o aumento del gasto público? Aclara por qué si /por qué no)}

Por tanto, este tipo de desgravación incentiva a la empresa a invertir hasta el punto en que la suma del valor del capital y de la devolución excede el coste del capital.

\textbf{Desgravación fiscal: Permanente}

\tetcolor{red}{\textbf{OJO}!} Entender bien la dinámica de ajuste -- No es fácil.

\textcolor{magenta}{
  desarrollar muy bien, como se dice: no es fácil. quiero entenderlo, ser claro, didáctico y pausado, no hay necesidad de comprimir explicaciones. Gracias.
}

Si la desgravación es permanente, $q^*_{EE}\equiv\dot K=0$ se desplaza hacia abajo, hacia el nuevo sendero de silla, en el mismo momento en que se produce la noticia. 

\imagenfit{ICEX_PF.png}{Efectos de una desgravación fiscal}{\textit{Macroeconomía avanzada}. ROMER, 2006}

Intuitivamente, la desgravación estimula la inversión no porque aumente la ratio de valor de capital instalado y coste de reposición, sino porque aminora el coste de adquisición del capital: una parte la paga el Estado. Esto nos permite resolver la aparente paradoja de aumento de la inversión y disminución de $q<1$. En este contexto es necesario $q_t>1-\gamma(t)$ para invertir, y no $q_t>1$. La secuencia sencilla es: (i) salto del nivel estacionario de inversión neta nula al nuevo nivel de $q_t$ dado el capital instalado $K_0^*$ que le sitúa sobre la senda de convergencia; (ii)  desplazamiento de $q$ y $K$ a lo largo del sendero de silla hacia el nuevo equilibrio a largo plazo; y, (iii) valor $K^*_1$ nuevo de equilibrio, $K_0^*<K_1^*$, y $q^*_{EE_1}=1-\gamma(t)$, menor que el inicial ($1$, punto $E'$).

\textbf{Desgravación fiscal ($\gamma$): Transitoria}

\textcolor{magenta}{Desarrollo analítico más profundo, qué pasa, por qué, implicaciones, etc.}

Los efectos de una desgravación temporal son muy diferente. En términos generales, observamos que el anuncio de una desgravación hace caer $q$ hasta un punto en que la evolución de $K$ y $q$, dada la desgravación, les conduce de nuevo hacia el sendero de silla inicial a medida que se agota el período de vigencia de la desgravación (el punto $B$ indica el fin de la vigencia de la desgravación):

\imagenfit{ICEX_PF_Transitoria.png}{Efectos de una desgravación fiscal transitoria}{\textit{Macroeconomía avanzada}. ROMER, 2006}

Esto es así porque, como se observa, $q$ no llega a alcanzar en este caso su valor sobre el nuevo sendero de silla (se sitúa sobre $A$); así, pues, una desgravación temporal reduce $q$ en una proporción menor: como la desgravación temporal no lleva a un incremento permanente del stock de capital, la disminución del valor del capital no es tan pronunciada como en la hipótesis anterior. 

Ahora bien, nótese que $q$ es mayor cuando la desgravación es transitoria que cuando es de naturaleza permanente, por tanto, una desgravación transitoria tiene un efecto mayor sobre la inversión que una permanente. 

Por último, puede observarse como en caso de una desgravación trasitorio, existe un momento en el que se produce un boom inversor ($B$). Es posible racionalizar esto situando este momento en el último periodo de vigencia de la desgravación: el ritmo inversor crece a medida que las empresas tratan de invertir justo antes de que la desgravación desaparezca. 

Vistos los diferentes efectos de política económica, cabe señalar que estos efectos serán diferentes dependiendo de si el cambio exógeno es permanente o temporal. Un cambio permanente da lugar a un salto inmediato al nuevo sendero de la silla, seguido de una convergencia a lo largo de la senda. Por el contrario, un cambio temporal no situará a la economía en la senda de equilibrio en tanto las variables exógenas cambiarán en el futuro (recuérdese que el sendero de silla describe la convergencia del sistema cuando los valores futuros de las variables exógenas son estables). Por tanto, la economía salta a un sendero de silla inestable que conduce de nuevo al sendero de silla original en el momento en que la política económica vuelva a su estado original.


\textbf{1.3.2.- Política Monetaria: Tipos de interés}

\textcolor{magenta}{
  Los efectos de una Política Monetaria son similares a los vistos en relación al aumento de la producción agregada, a excepción de los cambios en la pendiente de la curva $\dot q_t=0$.  (desarrollo analítico: ¿dónde incide? ¿ atravñes de qué canales?)
}

Partiendo de la ecuación de la curva $\dot q_t=rq_t-P_xF'(K_t)=0$, realizamos los siguientes cálculos:

$$\begin{aligned}
    &(1)\quad \dot q_t=rq_t-P_xF'(K_t)\quad\underset{\text{\scriptsize equilibrio}}{\Longrightarrow}\quad rq_t=P_xF'(K_t)\\[0.5em]
    &(2)\quad q_t=\frac{P_xF'(K_t)}{r}\\[0.5em]
    &(3)\quad q_t=\underbrace{\frac{P_x}{r}}_{\text{\scriptsize cte.}}F'(K_t)=\overset{\text{\scriptsize Senda de equilibrio}}{q(K_t)}\quad\underset{\text{\scriptsize Como, $F'(K_t)<0$}}{\Longrightarrow}\quad \frac{\mathrm d}{\mathrm dK}q(K_t)=\frac{P_x}{r}\frac{\mathrm d}{\mathrm dK}F'(K_t)\\[0.5em]
    &(4) \quad \boxed{\frac{\mathrm d}{\mathrm dK}q(K_t)=\frac{P_x}{r}F'(K_t}\quad\Longrightarrow\quad 
    \left\{\begin{aligned}
        &(1) &&\uparrow\frac{\mathrm d}{\mathrm dK_t}q(K_t)\\[0.5em]
        &(2) && \Delta (rq_t)=\Delta(P_xF'(K_t))\iff \dot q=0
  \end{aligned}\right.
\end{aligned}$$

\textcolor{magenta}{Desarrollar mucho mejor lo anterior, no se entiende muy bien, dividir y estructurar en párrafos, como se suele hacer siempre.}


\textbf{1.3.2.1.- Perturbaciones permanentes}

\textcolor{magenta}{En el caso de una caída permanente del tipo de interés, la curva $\dot q_t=0$ se desplaza hacia arriba (aumentando su pendiente), $q$ se desplaza al nuevo sendero de silla (punto $A$), mientras que $K$ y $q$ se desplazan hasta alcanzar el nuevo equilibrio a largo plazo (punto $E'$). (desarrollo analítico)}

\textbf{1.3.2.2.- Perturbaciones transitorios}

\textcolor{magenta}{(desarrollo analítico)}


\textbf{1.2.- La q marginal: ¿cómo se obtiene y qué relación tiene con la Q de TOBIN?}


\textcolor{magenta}{(Juntas lo anterior con la representación de la $q_t$ siguiente, teniendo en cuenta este párrafo del cuerpo) Como podemos ver, la variable $q_t$ sintetiza toda la información relevante sobre el futuro que la empresa necesita conocer para adoptar sus decisiones de inversión. ¿Por qué? Poruqe $q$ refleja los efectos que tendría una unidad adicional de capital sobre el valor presente de los beneficios futuros: es el valor presente descontado de los ingresos marginales futuros de una unidad adicional de capital.}

\textcolor{magenta}{Por tanto, podemos definir una formulación alternativa para este parámetro: el cociente entre el valor descontado de los beneficios futuros (valor unitario del capital instalado) en el cociente.

Por otra parte, la $q$ puede interpretarse como el valor de mercado de una unidad adicional: (¿cuál es la representación que sigue? Parece que se están diciendo dos cosas contradictorias, resolver esta cuestión, reescribir y desarrollar analíticamente la obtención de la fórmula siguiente de forma que se entrelace con la explicación verbal que se vaya a reescribir.)}

$$\begin{aligned}
  \text{Formulación marginal}:&&q=\frac{V'(I_t^o, K_t)}{p_k C'(I_t^o,K_t)}\quad 
\end{aligned}$$

Asumiremos que $\bar p_k=1$, es decir, \textcolor{magenta}{explicar qué significa esta suposición y por qué es válida (entiendo que este supuesto es donde se incorpora el supuesto de : no existen costes externos; ¿Podría eliminarlo? ¿Podríamos continuar considerando que sí existen? ¿cambiaría mucho la formulación?}. Incorporando este supuesto a la fórmula anterior, acabamos obteniendo:


$$\begin{aligned}
  \quad \boxed{q=\frac{V'(I_t^o, K_t)}{C'(I_t^o,K_t)}}\quad\Longrightarrow
    \left\{\begin{aligned}
        &q>1\equiv\text{Inversión rentable}\\[0.5em]
        &q<1\equiv\text{Inversión no rentable}\\[0.5em]
    \end{aligned}\right.
\end{aligned}$$

Es decir, el \textbf{cociente de optimalidad marginal}: valor marginal de capital instalado entre el valor marginal del capital de nueva adquisición. \textcolor{magenta}{¿Qué implicaciones tiene? ¿qué es?...}


\begin{specialblock}
  Como sabemos, la $q$ de TOBIN puede considerarse como el cociente entre el valor de mercado del capital instalado y el coste de reposición:
  
  $$\begin{aligned}
      &\text{Q de TOBIN}:&& {q=\frac{\text{Valor de mercado}}{\text{Coste de reposición}}}\Longrightarrow
      \left\{\begin{aligned}
          &q>1\equiv\text{Rentable}\\[0.5em]
          &q<1\equiv\text{No rentable}
      \end{aligned}\right.
  \end{aligned}$$
  
  Es decir, la relación entre el \textbf{valor total de la empresa} y el \textbf{coste de reposición de su stock de capital total}.

  --------\textcolor{magenta}{De todo lo que está entre ----- debes complementar el razonamiento con herramientas matemáticas. Es decir, explicar lo que estás diciendo en palabras con su representación matemática. Además, deberíamos explicar que pasaría si no se cumpliera el supuesto: ¿qué resultados obtendriamos? ¿darían lecturas/recomendaciones/implicaciones complementarias o contradictorias?¿por qué? Todo esto es importante que se explique bien y se vea acompañado con la representación matemática}

  Ahora bien, por qué la equivalencia con lo anterior depende, estríctamente, de la presencia de rendimientos de escala constantes en los costes de ajuste (HAYASHI, 1982). 
  
  Para entender por qué, hay que recordar que la empresa se enfrenta a costes de ajuste convexos. 

  Por tanto, que haya rendimientos constantes en los costes de ajuste significa algo muy concreto: si duplicamos el tamaño de la empresa (es decir, duplicamos $K$ e $I^o$), el coste de ajuste también se duplica. 
  
  Formalmente, la función es homogénea de grado $1$: $C(\lambda I, \lambda K) = \lambda C(I,K)$. Esto tiene una implicación crucial: \textbf{el coste de ajuste no depende del tamaño absoluto de la empresa} ($K_t$), sino de la inversión relativa, es decir, de la ratio $I/K$, que es \textbf{la tasa de crecimiento del capital}. 
  
  Por tanto, el problema dinámico de la empresa se puede normalizar por $K$ y la variable de control relevante pasa a ser únicamente la tasa de crecimiento del capital.

  La $q$ marginal, entonces, deja de depender del nivel de $K$ y pasa a ser una función únicamente de esa tasa de crecimiento. Y, al mismo tiempo, la $q$ media también se convierte en una magnitud independiente de la escala. Por eso ambas coinciden: porque ambas están describiendo lo mismo en un modelo donde todo es proporcional.

  En definitiva, cuando hay rendimientos constantes en los costes de ajuste, no hay efectos de tamaño. Una empresa grande y una pequeña enfrentan exactamente el mismo problema en términos relativos. \textbf{Lo único que importa es cuánto crecen, no cuánto tienen}. 
  
  Una consecuencia directa, como veremos, es que todas las empresas optan por la misma tasa de crecimiento de su stock de capital y de su $q$ de TOBIN: \textbf{no hay convergencia ni divergencia}.
  
  -------------
\end{specialblock}

\clearpage
\anexo{Derivación: Modelos de información imperfecta}\label{anx:modelos_imperfecta}

La mayoría de las decisiones de inversión presentan en la realidad una serie de rasgos comunes:
\begin{lnum}
    \item \textbf{Incertidumbre permanente}. En primer lugar, el contexto económico tiene una incertidumbre permanente y la información llega de forma gradual; y,
    \item \textbf{Información asimétrica}. En segundo lugar, la información, que llega de forma gradual, no es la misma para todos los participantes. En particular, varios competidores en una misma industria tendrán una información completa sobre sí mismos, una información incompleta sobre el mercado, y un grado de asimetría en la información que disponen de entre sí;
\end{lnum}   

\textbf{1.- Riesgo: ¿cómo afecta el riesgo a las decisiones de inversión?}






\textbf{1.1.- Incertidumbre permanente: ¿cómo afectan las previsiones a las decisiones de inversión?}

Una forma razonable de introducir un contexto de riesgo es mediante HER: las empresas cuentan con buenos equipos, pueden contratar servicios externos, disponer de información en tiempo real, y tomar sus decisiones de inversión futuras en base previsiones fundadas que todo este conjunto de agentes y medios les proporciona.

Introduzcamos, pues, HER en el problema de optimización anterior:

\textcolor{magenta}{(Hay que hacer todo el desarrollo analítico con HER) Además, debe irse intercalando con desarrolo escrito, tal y como se hizo en la parte anterior.}




\textcolor{magenta}{(Aquí terminaba el desarrollo anterior, ver si termina aquí o necesita más pasos)}Resolviendolo, obtendremos la siguiente CPO:

\eqblock{
\begin{aligned}
    \text{HER}\Longrightarrow\quad \boxed{\mathbb E_t[\dot q_t]=rq_t-PF'(K_t)}
\end{aligned}
}{}

Por tanto, cada empresa invierte hasta el punto en que el coste de adquisición del nuevo capital iguala su valor de mercado: la ecuación relativa a la optimalidad marginal de la inversión sigue siendo válida (primera CPO, $q_t=1+C'(I_t^o)$).


\textbf{1.2.- Irreversibilidad: Costes de ajuste asimétricos}

Otra forma en la que se manifiesta la información incompleta es mediante las funciones de costes de ajuste asimétricos.

En la mayoría de situaciones, la empresa debe tener presentes los costes que afrontaría ante un escenario desfavorable y, en general, esos costes están muy sesgados hacia abajo: es más caro desinvertir que invertir. 

Paa incluir este hecho en nuestra problema, debemos realizar una serie de moficiaciones básicas.

\textcolor{magenta}{(incorporar el desarrollo analítico de este problema). Ademá,s muy importante, ir intercalando con párrafo de explicación, igual que siempre.

redesarrollar todo lo de abajo:}

Ante esta situación, la empresa estaría experimentando costes hundidos a la hora de tomar sus decisiones y debería comportarse de manera más aversa al riesgo respecto a las contingencias futuros, puesto que instalar capacidad es menos costoso que deshacerse de ello ante un cambio de las circunstancias:


$$\begin{aligned}
    \text{HER + Costes hundidos}\Longrightarrow\text{Aversión}
    \left\{\begin{aligned}
        &C(\cdot)\not\in\text{Poliniomios}\Longrightarrow\underset{\text{\scriptsize p.e. lineal $\in\mathbb R^+$ y cuadrática $\in\mathbb R^-$}}{\text{Asimétrica}}\\[0.5em]
        &\mathbb E_t[\dot q_t]=rq_t-P_xF'(K_t)
    \end{aligned}\right.
\end{aligned}$$

En este contexto, si $q(K_t)$ (o $K(q_t)$, ya que $q(K_t)=K^{-1}(q_t)$) se desplaza hacia arriba, el stock de capital total de la industria aumenta rápidamente ($\dot q_t\gg0$); pero si se desplaza hacia abajo, $K$ se reduce gradualmente ($\dot q_t\sim0$). Si analizamos el ejemplo anterior en presencia de costes de ajuste asimétricos, cuando la pretensión del gobierno de modificar la política fiscal se hace pública, $q_t$ salta hacia arriba hasta un punto en que $K$ y $q$ se sitúan en la senda estable (del momento de la votación). 

\imagenfit{ICEX_Irreversibilidad.png}{Irreversibilidad en la función de inversión}{\textit{Macroeconomía avanzada}. ROMER, 2006}

Sin embargo, el salto es en esta ocasión menor que el que tenía lugar cuando presuponíamos unos costes de ajuste simétricos. Esto se debe a que el sendero de silla es curvo: si el valor de $K$ excede a su valor de equilibrio a largo plazo, $K$ disminuye lentamente y $q$ es mucho menor que $1$ (pero se ajusta muy lentamente, $\dot q_t\sim 0$). Por el contrario, si $K$ se encuentra por debajo de su valor de equilibrio, a largo plazo aumenta rápidamente, de manera que $q$ es sólo ligeramente superior a la unidad (pero se ajusta muy rápidamente, $\dot q_t\gg0$). 

¿Por qué? Porque si las empresas acumulan un stock de capital importante antes de la votación y la propuesta resulta luego rechazada, la presencia de costes de ajuste dificultaría el proceso de desinversión. Esto contribuye a reducir el valor del capital antes de la votación y, por consiguiente, la inversión. En definitiva, si los costes de ajuste son asimétricos, la incertidumbre sobre la posición de la función de beneficios reduce las expectativas sobre rentabilidad futura, y con ellas la demanda de inversión. Esta asimetría implica que la inversión es en cierto sentido \textbf{irreversible}: aumentar el stock de capital es más sencillo que revertir esta decisión.\footnote{%
    Por ejemplo, en el caso de que una empresa construya una nueva planta para aumentar su capacidad de exportación con un tamaño específico y una tecnología específica, la empresa se vería limitada en sus posibilidades de decisión si en un futuro desease cambiar el tamaño de la planta.
} 

De forma intuitiva, cuando la inversión es irreversible, existe un valor de opción asociado a la espera: si la empresa no invierte, conserva al menos la posibilidad de poseer un stock de capital bajo, mientras que, si invierte, se compromete a retener un stock de capital elevado.

Además, conviene señalar que la irreversibilidad no es un problema exclusivamente microeconómico ya que no se diluye con la agregación en tanto suponemos empresas heterogéneas. Por ejemplo, la incertidumbre sobre las variables macroeconómicas (tipos de interés, tipos de cambio, tipos impositivos, etc.) afectaría a todas las empresas.\footnote{%
    Bernanke (1983) considera que el nivel de incertidumbre percibido por las empresas varía en función del ciclo económico y enfatiza la importancia de la irreversibilidad para analizar el comportamiento cíclico de la inversión agregada.
}


\textbf{1.3.- Opciones reales: McDONALD y SIEGEL (1986)}

A continuación, analizaremos el enfoque de las opciones reales propuesto por McDonald y Siegel (1986). Estos autores analizan un proyecto de tamaño fijo, por lo que el calendario del proyecto es la única opción que hay que elegir.\footnote{%
  Los trabajos posteriores de PINDYCK (1988) y BERTOLA (1988), permiten explicar la inversión incremental, de modo que que la empresa elige tanto el momento como el tamaño de sus inversiones.
} 

En su modelo la inversión irreversible se considera una opción real, es decir, un derecho y no una obligación de llevar a cabo una inversión en capital físico. De esta forma cuando una empresa lleva a cabo un gasto de inversión irreversible ejerce su opción de invertir, renunciando a la posibilidad de esperar para obtener nueva información que podría afectar al nivel de capital productivo instalado. 

Esta opción perdida se considera un coste de oportunidad y se incluye, por tanto, como parte del coste de la inversión. Para tener en cuenta este aspecto deber modificarse la regla de decisión para invertir: el rendimiento esperado debe ser mayor que el coste de adquisición e instalación en la cuantía del valor de mantener la opción viva:

\textcolor{magenta}{(incluir completo el desarrollo analítico, complementando desarrollo con párrafos que aclaran las operaciones) Todo lo de abajo hasta ----}

$$\begin{aligned}
    \text{Optimalidad del proyecto}:q=\frac{V}{K}\quad \underset{\text{\scriptsize con opción real}}{\Longrightarrow}\quad
    \left\{\begin{aligned}
        (\checkmark)\iff q>1+p_{\text{\scriptsize opción}}\\[0.5em]
        (\mathrm X) \iff q\leq1+p_{\text{\scriptsize opción}}
    \end{aligned}\right.\\[1em]
    \Longrightarrow
    \begin{aligned}
        &(1)&&\text{Retardos óptimos}\\[0.5em]
        &(2)&&\text{Acentúa la incertidumbre}: \frac{\partial p_{\text{opción}}}{\partial \text{Risk}}>0\\[0.5em]
        &(3)&&\underset{\substack{\text{\scriptsize posponer es óptimo}\\\text{\scriptsize por la información}}}{\text{Histéresis}}:\frac{\partial p_{\text{opción}}}{\partial t}<0
    \end{aligned}
\end{aligned}$$

Esta modificación de la regla de decisión de inversión tiene interesantes implicaciones.

En primer lugar, implica que los \textbf{retardos son óptimos}: la presencia de esta opción implica que es óptimo retrasar la inversión, en lugar de llevarla a cabo inmediatamente, incluso cuando la ejecución inmediata tenga un valor positivo. En cambio, el valor de la inversión puede incrementarse si se espera para obtener información adicional.

En segundo lugar, \textbf{acentúa el papel de la incertidumbre}: como la mayoría de las opciones, el valor de esperar ($p_{\text{\tiny opción}}$) es creciente con la incertidumbre. Esta característica implica un efecto de la incertidumbre sobre el valor y el calendario o momento de ejecución de las inversiones que no existe en la mayoría de los modelos convencionales.

En tercer lugar, el fenómeno de \textbf{histéresis}: DIXIT (1992) habla de modelos de inercia óptima o tiranía benevolente del status quo, al centrarse en los fenómenos de histéresis (persistencia de un efecto cuando ya ha desaparecido su causa). En concreto, este autor señala que en un contexto de incertidumbre se genera un trade-off entre las ventajas de invertir cuanto antes y las de esperar la llegada de nueva información, de forma que las empresas posponen sus decisiones de inversión para mantener sus opciones abiertas. En otras palabras, las empresas se niegan a invertir incluso cuando el rendimiento esperado en la actualidad supera con el coste del capital, ya que pueden optar por esperar para estar más seguros de que esa situación favorable no es transitoria.

A título ilustrativo, este fenómeno de histéresis de inversión se observa en la lenta respuesta de las importaciones estadounidenses a las variaciones del tipo de cambio durante los años 80. Desde 1980 hasta finales de 1984, el dólar estadounidense se apreció aproximadamente un 50\%, de forma que la ventaja competitiva de las empresas extranjeras en el mercado estadounidense aumentó drásticamente. Sin embargo, el volumen de las importaciones no empezó a aumentar hasta principios de 1983, lo que supone un retraso mucho mayor que el año o los 18 meses que se consideraban hechos estilizados hasta entonces. En 1985 el dólar empezó a depreciarse y a finales de 1987 casi había recuperado su nivel de 1978, pero el volumen de las importaciones no disminuyó durante otros dos años (sufriendo incluso ligeros aumentos). Esta histéresis en la inversión puede explicar en tanto, una vez establecidas en el mercado estadounidense, las empresas extranjeras actuaron con lentitud a la hora de reducir la escala de sus exportaciones cuando el tipo de cambio comenzó a ser desfavorable. 

------

La investigación sobre la irreversibilidad de la inversión se ha ramificado tanto empírica como teóricamente. En relación a los estudios empíricos, las primeras aplicaciones se centraban en los recursos naturales (Brennan y Schwartz, 1985), mientras que posteriormente los estudios han incluido los bienes inmuebles o la inversión en equipos (Eberly y van bienes de capital (Eberly y van Mieghem, 1997). En cuanto a los estudios teóricos, la agregación de modelos con ajustes poco frecuentes sigue siendo un reto y los resultados siguen siendo controvertidos. Salvo en casos muy especiales (Caplin y Spulber, 1987), la agregación requiere el seguimiento de una distribución de agentes. Sin embargo, es precisamente esta característica la que puede coincidir con la observación de que gran parte de la volatilidad en la inversión en la práctica surge del margen extensivo (el número de agentes que se ajustan) y no del margen intensivo (el tamaño medio del ajuste).



\textbf{1.4.- Información y paciencia: BERNANKE}

\textcolor{magenta}{Otra forma de concebir la irreversibilidad es la introducida por BERNANKE: la función de inversión es muy sensible a la información. 
    
De esta forma, el autor plantea un un trade off entre las ventajas de invertir cuanto antes (adquirir ventajas competitivas como reputación o disfrute de patentes) y las de esperar a la llegada de nueva información. 
    
Esto da lugar a un proceso de decisión de carácter discontinuo, en el cual las empresas alternan periodos de inversión con periodos de inacción durante los cuales se permite la depreciación del capital existente. Así, al igual que otras formas de costes de ajuste, la irreversibilidad en las decisiones de inversión reduce la sensibilidad de la demanda de inversión ante shocks exógenos. 

Así, la medida en la que la inacción a nivel empresa afecta a la dinámica agregada depende del grado de sincronización de las acciones individuales. (redesarrollar y explicar mejor, no lo entiendo muy bien. Gracias.)}
    

----------



A título de ejemplo BLOOM, BOND, y VAN REENEN (2007) estudian el comportamiento de las inversiones de una muestra de 672 empresas manufactureras del Reino Unido que cotizan en bolsa durante el periodo 1972-1991 e indican que incrementos en la incertidumbre durante grandes recesiones, como la crisis del petróleo de los años 70s, podrían reducir considerablemente la respuesta de la inversión ante Políticas Monetarias o Fiscales llevadas a cabo por las autoridades. En términos generales, los autores observan una persistencia de esta inacción en la inversión empresarial ante periodos de alta incertidumbre. 
    
No obstante, puede haber otras justificaciones que podrían explicar las mismas pautas en la dinámica de inversión de las empresas, por ejemplo, que las empresas sometidas a una mayor incertidumbre otorguen menos importancia a la información reciente a la hora de actualizar sus expectativas de crecimiento futuro. 

\textbf{2.- Información asimétrica: ¿?}



\textbf{Mercados de capital: ¿?}


--------
\textcolor{magenta}{(Desarrollar el modelo que se presenta a continuación) Todo entero, analíticamente y complementado, como siempre. Gracias.}

La manera más sencilla de concebir esta estructura no lineal de los costes de financiación es a través de los modelos Principal-Agente: cuando existe información asimétrica, la inversión no sólo depende de los tipos de  interés y de la rentabilidad, sino de factores adicionales, como la capacidad del inversor para controlar la empresa o la capacidad de ésta para financiar la inversión mediante recursos internos.\footnote{%
    Principalmente, este comportamiento no lineal de los costes viene ocasionado por lo siguiente:
\begin{lnum}
    \item \textbf{Costes deerivados de la asimetría informativa que enfrenta el Principal}. De este modo una parte del riesgo asociado a la rentabilidad de la inversión lo asumen los inversores. Esto es lo que ocurre, por ejemplo, cuando existe la posibilidad de que la empresa quiebre. Si esto sucede, la empresa puede modificar su comportamiento para aprovechar su ventaja relativa en materia de información: puede decidir financiar la inversión exclusivamente mediante préstamos u optar por una estrategia de alto riesgo en lugar de asumir una de bajo riesgo aun a sabiendas de que esto reduce la rentabilidad esperada del proyecto. 
    \item \textbf{Costes derivados del comportamiento del agente}. Otro tipo de imperfección es aquella derivada de los costes de agencia. Téngase en cuenta que, si el nivel de endeudamiento es muy elevado, y dada la existencia de responsabilidad limitada por parte de los accionistas, éstos pueden estar interesados en proyectos rentables pero arriesgados, que podrían poner en peligro la solvencia de la empresa. 
\end{lnum}
Ante este comportamiento, los prestamistas impondrán una serie de restricciones que, en general, tenderán a elevar el coste de los préstamos.
}

Este tipo de imperfecciones que condicionan las decisiones de financiación de la empresa ha dado lugar a la Teoría de la jerarquía financiera, a través de la cual distintas fuentes de financiación se ordenan según su coste: la de menor coste resulta ser la financiación con beneficios retenidos, después está el endeudamiento y, finalmente, la emisión de nuevas acciones. Asimismo, se supone que el coste de la deuda es creciente con el nivel de deuda requerida.  El análisis de las consecuencias de esta teoría para las decisiones de inversión permite establecer una tipología de empresas susceptible de representación gráfica (BOND y MEGHIR, 1990):

\imagenfit{ICES_TJF.png}{Teoría de la Jerarquía Financiera}{\textit{Macroeconomía Avanzada I}. ARGANDOÑA, A.; GÁMEZ, C. y MOCHÓN, F. (1997)}

Por un lado, estarán las empresas que generan recursos (beneficios y dotación a amortizaciones) suficientes para cubrir sus oportunidades de inversión sin recurrir a la financiación externa (situación $I_1$). Un segundo grupo de empresas son aquellas que no generan recursos suficientes y recurren a la emisión de deuda, cuyo coste será creciente conforme aumente el nivel requerido de recursos (situación $I_2$). Una tercera categoría viene determinada cuando el coste de la deuda adicional es tal que a la empresa le interesa no endeudarse más y recurrir a una ampliación de capital emitiendo nuevas acciones, con un coste fijo (situación $I_3$).\footnote{%
    Téngase en cuenta que las empresas del grupo $I_1$ disponen de fondos internos suficientes para financiar nuevos proyectos de inversión sin que aumente el coste marginal de la financiación. Por ello, para este tipo de empresas, un aumento exógeno de los recursos internos no incide sobre la inversión, trasladándose íntegramente a dividendos, lo que es una muestra de que disponen de holgura financiera. No obstante, en el caso de las empresas del grupo $I_2$ el volumen de inversión sí se verá afectado por el aumento de los recursos internos, pues se enfrentan a un coste marginal de la inversión creciente. Por ello diremos que las empresas que están en esta situación (que no reparten dividendos ni emiten nuevas acciones) no disponen de holgura financiera.
}

------



\end{document}